\documentclass[11pt,fleqn, openany]{book} % Default font size and left-justified equations

%%%%%%%%%%%%%%%%%%%%%%%%%%%%%%%%%%%%%%%%%
% The Legrand Orange Book
% Structural Definitions File
% Version 2.1 (26/09/2018)
%
% Original author:
% Mathias Legrand (legrand.mathias@gmail.com) with modifications by:
% Vel (vel@latextemplates.com)
% 
% This file was downloaded from:
% http://www.LaTeXTemplates.com
%
% License:
% CC BY-NC-SA 3.0 (http://creativecommons.org/licenses/by-nc-sa/3.0/)
%
%%%%%%%%%%%%%%%%%%%%%%%%%%%%%%%%%%%%%%%%%

%----------------------------------------------------------------------------------------
%	VARIOUS REQUIRED PACKAGES AND CONFIGURATIONS
%----------------------------------------------------------------------------------------

\usepackage[table]{xcolor}

\usepackage{graphicx}
\usepackage{tabularx} % Required for including pictures
\usepackage{pgf,tikz,tkz-tab,eurosym,yhmath, stmaryrd}
\usepackage{pgfplots}
\usepackage{mathrsfs}
\usetikzlibrary{patterns}
\usetikzlibrary{trees}
\graphicspath{{../../Pictures/}}
\usepackage{multicol} 


\usepackage[english]{babel} % English language/hyphenation
\usepackage{icomma}
\usepackage{enumitem} % Customize lists
\setlist{nolistsep, nosep, nolistsep} % Reduce spacing between bullet points and numbered lists

\usepackage{booktabs} % Required for nicer horizontal rules in tables

 % Required for specifying colors by name


\definecolor{ocre}{RGB}{243,102,25} % Define the orange color used for highlighting throughout the book

\usepackage{listings}

\definecolor{codegreen}{rgb}{0,0.6,0}
\definecolor{codegray}{rgb}{0.5,0.5,0.5}
\definecolor{codepurple}{rgb}{0.58,0,0.82}
\definecolor{backcolour}{rgb}{0.95,0.95,0.92}

\lstdefinestyle{mystyle}{
    backgroundcolor=\color{backcolour},   
    commentstyle=\color{codegreen},
    keywordstyle=\color{magenta},
    numberstyle=\tiny\color{codegray},
    stringstyle=\color{codepurple},
    basicstyle=\ttfamily\footnotesize,
    breakatwhitespace=false,         
    breaklines=true,                 
    captionpos=b,                    
    keepspaces=true,                 
    numbers=left,                    
    numbersep=5pt,                  
    showspaces=false,                
    showstringspaces=false,
    showtabs=false,                  
    tabsize=2
}

\lstset{style=mystyle}

%----------------------------------------------------------------------------------------
% Paramétrage XSIM
%----------------------------------------------------------------------------------------

\usepackage[no-files]{xsim}


\DeclareExerciseEnvironmentTemplate{myex}{%
    \textbf{%
      \hypertarget{ex:\ExerciseID}{\sffamily{\ensuremath{\blacktriangleright}} Exercice \GetExerciseProperty{counter} \GetExerciseProperty{subtitle} --}
      \hyperlink{sol:\ExerciseID}{Voir le corrigé}%
    }\par
}{\par\smallskip}

\DeclareExerciseEnvironmentTemplate{mysol}{%
    \textbf{%
      \hypertarget{sol:\ExerciseID}{\sffamily{\ensuremath{\blacktriangleright}} Correction \GetExerciseProperty{counter} --}
      \hyperlink{ex:\ExerciseID}{Voir l'énoncé}%
    }\par
}{\par\medskip}

\xsimsetup{
  exercise/template = myex ,
  solution/template = mysol 
}

%Collection exercices

\DeclareExerciseTagging{topic}

\xsimsetup{collect}

%----------------------------------------------------------------------------------------
% SYMBOLES
%----------------------------------------------------------------------------------------

\newcommand\imCMsym[4][\mathord]{%
  \DeclareFontFamily{U} {#2}{}
  \DeclareFontShape{U}{#2}{m}{n}{
    <-6> #25
    <6-7> #26
    <7-8> #27
    <8-9> #28
    <9-10> #29
    <10-12> #210
    <12-> #212}{}
  \DeclareSymbolFont{CM#2} {U} {#2}{m}{n}
  \DeclareMathSymbol{#4}{#1}{CM#2}{#3}
}
\newcommand\alsoimCMsym[4][\mathord]{\DeclareMathSymbol{#4}{#1}{CM#2}{#3}}

\imCMsym{cmmi}{124}{\CMjmath}

\newcommand{\Oij}{(O\,;\,\vec{\imath}\,,\, \vec{\CMjmath} )}
\newcommand{\Oijk}{(O\,;\,\vec{\imath}\,,\, \vec{\CMjmath}\,,\,\vec{k})}

\newcommand\e{\mathrm{e}}
\newcommand\R{\mathbb{R}}
\newcommand\N{\mathbb{N}}


%----------------------------------------------------------------------------------------
%	MARGINS
%----------------------------------------------------------------------------------------

\usepackage{geometry} % Required for adjusting page dimensions and margins

\geometry{
	paper=a4paper, % Paper size, change to letterpaper for US letter size
	top=3cm, % Top margin
	bottom=3cm, % Bottom margin
	left=2cm, % Left margin
	right=2cm, % Right margin
	headheight=14pt, % Header height
	footskip=1.4cm, % Space from the bottom margin to the baseline of the footer
	headsep=10pt, % Space from the top margin to the baseline of the header
	%showframe, % Uncomment to show how the type block is set on the page
}

\setlength{\parindent}{0pt}
\parskip=5pt



%----------------------------------------------------------------------------------------
%	FONTS
%----------------------------------------------------------------------------------------

\usepackage{avant} % Use the Avantgarde font for headings
\usepackage{times} % Use the Times font for headings
\usepackage{mathptmx} % Use the Adobe Times Roman as the default text font together with math symbols from the Sym­bol, Chancery and Com­puter Modern fonts

%\usepackage{microtype} % Slightly tweak font spacing for aesthetics
%\usepackage[utf8]{inputenc} % Required for including letters with accents
\usepackage[T1]{fontenc} % Use 8-bit encoding that has 256 glyphs

%----------------------------------------------------------------------------------------
%	BIBLIOGRAPHY AND INDEX
%----------------------------------------------------------------------------------------

\usepackage[style=numeric,citestyle=numeric,sorting=nyt,sortcites=true,autopunct=true,babel=hyphen,hyperref=true,abbreviate=false,backref=true,backend=biber]{biblatex}
\addbibresource{bibliography.bib} % BibTeX bibliography file
\defbibheading{bibempty}{}

\usepackage{calc} % For simpler calculation - used for spacing the index letter headings correctly
\usepackage{makeidx} % Required to make an index
\makeindex % Tells LaTeX to create the files required for indexing

%----------------------------------------------------------------------------------------
%	MAIN TABLE OF CONTENTS
%----------------------------------------------------------------------------------------

\usepackage{titletoc} % Required for manipulating the table of contents

\contentsmargin{0cm} % Removes the default margin

% Part text styling (this is mostly taken care of in the PART HEADINGS section of this file)
\titlecontents{part}
	[0cm] % Left indentation
	{\addvspace{20pt}\bfseries} % Spacing and font options for parts
	{}
	{}
	{}

% Chapter text styling
\titlecontents{chapter}
	[1.25cm] % Left indentation
	{\addvspace{12pt}\large\sffamily\bfseries} % Spacing and font options for chapters
	{\color{ocre!60}\contentslabel[\Large\thecontentslabel]{1.25cm}\color{ocre}} % Formatting of numbered sections of this type
	{\color{ocre}} % Formatting of numberless sections of this type
	{\color{ocre!60}\normalsize\;\titlerule*[.5pc]{.}\;\thecontentspage} % Formatting of the filler to the right of the heading and the page number

% Section text styling
\titlecontents{section}
	[1.25cm] % Left indentation
	{\addvspace{3pt}\sffamily\bfseries} % Spacing and font options for sections
	{\contentslabel[\thecontentslabel]{1.25cm}} % Formatting of numbered sections of this type
	{} % Formatting of numberless sections of this type
	{\hfill\color{black}\thecontentspage} % Formatting of the filler to the right of the heading and the page number

% Subsection text styling
\titlecontents{subsection}
	[1.25cm] % Left indentation
	{\addvspace{1pt}\sffamily\small} % Spacing and font options for subsections
	{\contentslabel[\thecontentslabel]{1.25cm}} % Formatting of numbered sections of this type
	{} % Formatting of numberless sections of this type
	{\ \titlerule*[.5pc]{.}\;\thecontentspage} % Formatting of the filler to the right of the heading and the page number

% Figure text styling
\titlecontents{figure}
	[1.25cm] % Left indentation
	{\addvspace{1pt}\sffamily\small} % Spacing and font options for figures
	{\thecontentslabel\hspace*{1em}} % Formatting of numbered sections of this type
	{} % Formatting of numberless sections of this type
	{\ \titlerule*[.5pc]{.}\;\thecontentspage} % Formatting of the filler to the right of the heading and the page number

% Table text styling
\titlecontents{table}
	[1.25cm] % Left indentation
	{\addvspace{1pt}\sffamily\small} % Spacing and font options for tables
	{\thecontentslabel\hspace*{1em}} % Formatting of numbered sections of this type
	{} % Formatting of numberless sections of this type
	{\ \titlerule*[.5pc]{.}\;\thecontentspage} % Formatting of the filler to the right of the heading and the page number

%----------------------------------------------------------------------------------------
%	MINI TABLE OF CONTENTS IN PART HEADS
%----------------------------------------------------------------------------------------

% Chapter text styling
\titlecontents{lchapter}
	[0em] % Left indentation
	{\addvspace{15pt}\large\sffamily\bfseries} % Spacing and font options for chapters
	{\color{ocre}\contentslabel[\Large\thecontentslabel]{1.25cm}\color{ocre}} % Chapter number
	{}  
	{\color{ocre}\normalsize\sffamily\bfseries\;\titlerule*[.5pc]{.}\;\thecontentspage} % Page number

% Section text styling
\titlecontents{lsection}
	[0em] % Left indentation
	{\sffamily\small} % Spacing and font options for sections
	{\contentslabel[\thecontentslabel]{1.25cm}} % Section number
	{}
	{}

% Subsection text styling (note these aren't shown by default, display them by searchings this file for tocdepth and reading the commented text)
\titlecontents{lsubsection}
	[.5em] % Left indentation
	{\sffamily\footnotesize} % Spacing and font options for subsections
	{\contentslabel[\thecontentslabel]{1.25cm}}
	{}
	{}

%----------------------------------------------------------------------------------------
%	HEADERS AND FOOTERS
%----------------------------------------------------------------------------------------


\usepackage{fancyhdr} % Required for header and footer configuration

\pagestyle{fancy}
\renewcommand{\chaptermark}[1]{\markboth{\sffamily\normalsize\bfseries\ \thechapter.\ #1}{}} % Chapter text font settings
\renewcommand{\sectionmark}[1]{\markright{\sffamily\normalsize\thesection\hspace{5pt}#1}{}} % Section text font settings
\fancyhf{} \fancyhead[LE,RO]{\sffamily\normalsize\thepage} % Font setting for the page number in the header
\fancyhead[LO]{\rightmark} % Print the nearest section name on the left side of odd pages
\fancyhead[RE]{\leftmark} % Print the current chapter name on the right side of even pages

\fancyfoot[L]{Jason LAPEYRONNIE}
\fancyfoot[R]{\href{http://mathoutils.fr}{http://mathoutils.fr}} % Uncomment to include a footer

\renewcommand{\headrulewidth}{0.5pt} % Thickness of the rule under the header
\renewcommand{\footrulewidth}{0.5pt} % Thickness of the rule under the header

\fancypagestyle{plain}{% Style for when a plain pagestyle is specified
	\fancyhead{}\renewcommand{\headrulewidth}{0pt}%
}

% Removes the header from odd empty pages at the end of chapters
\makeatletter
\renewcommand{\cleardoublepage}{
\clearpage\ifodd\c@page\else
\hbox{}
\vspace*{\fill}
\thispagestyle{empty}
\newpage
\fi}

%----------------------------------------------------------------------------------------
%	THEOREM STYLES
%----------------------------------------------------------------------------------------

\usepackage{amsmath,amsfonts,amssymb,amsthm} % For math equations, theorems, symbols, etc

\newcommand{\intoo}[2]{\mathopen{]}#1\,;#2\mathclose{[}}
\newcommand{\ud}{\mathop{\mathrm{{}d}}\mathopen{}}
\newcommand{\intff}[2]{\mathopen{[}#1\,;#2\mathclose{]}}
\renewcommand{\qedsymbol}{$\blacksquare$}
\newtheorem{notation}{Notation}[section]

% Boxed/framed environments
\newtheoremstyle{ocrenumbox}% Theorem style name
{0pt}% Space above
{0pt}% Space below
{\normalfont}% Body font
{}% Indent amount
{\small\bf\sffamily\color{ocre}}% Theorem head font
{\;:\;}% Punctuation after theorem head
{0.25em}% Space after theorem head
{\small\sffamily\color{ocre}\thmname{#1}\nobreakspace\thmnumber{\@ifnotempty{#1}{}\@upn{#2}}% Theorem text (e.g. Theorem 2.1)
\thmnote{\nobreakspace\the\thm@notefont\sffamily\bfseries\color{black}---\nobreakspace#3}} % Optional theorem note

\newtheoremstyle{blacknumex}% Theorem style name
{5pt}% Space above
{10pt}% Space below
{\normalfont}% Body font
{} % Indent amount
{\small\bf\sffamily}% Theorem head font
{\;:\;}% Punctuation after theorem head
{0.25em}% Space after theorem head
{\small\sffamily{\tiny\ensuremath{\blacksquare}}\nobreakspace\thmname{#1}\nobreakspace\thmnumber{\@ifnotempty{#1}{}\@upn{#2}}% Theorem text (e.g. Theorem 2.1)
\thmnote{\nobreakspace\the\thm@notefont\sffamily\bfseries---\nobreakspace#3}}% Optional theorem note

\newtheoremstyle{blacknumexo}% Theorem style name
{15pt}% Space above
{10pt}% Space below
{\normalfont}% Body font
{} % Indent amount
{\small\bf\sffamily}% Theorem head font
{}% Punctuation after theorem head
{0.5em}% Space after theorem head
{\small\sffamily{\ensuremath{\blacktriangleright}}\nobreakspace\thmname{#1}\nobreakspace\thmnumber{\@ifnotempty{#1}{}\@upn{#2}}% Theorem text (e.g. Theorem 2.1)
\thmnote{\nobreakspace\the\thm@notefont\sffamily\bfseries---\nobreakspace#3} \\}% Optional theorem note



\newtheoremstyle{blacknumbox} % Theorem style name
{0pt}% Space above
{5pt}% Space below
{}% Body font
{}% Indent amount
{\large\bf\sffamily}% Theorem head font
{\;:\;}% Punctuation after theorem head
{0.25em}% Space after theorem head
{\small\sffamily\thmname{#1}\nobreakspace\thmnumber{\@ifnotempty{#1}{}\@upn{#2}}% Theorem text (e.g. Theorem 2.1)
\thmnote{\nobreakspace\the\thm@notefont\sffamily\bfseries---\nobreakspace#3}}% Optional theorem note

% Non-boxed/non-framed environments
\newtheoremstyle{ocrenum}% Theorem style name
{5pt}% Space above
{5pt}% Space below
{\normalfont}% Body font
{}% Indent amount
{\small\bf\sffamily\color{ocre}}% Theorem head font
{\;:\;}% Punctuation after theorem head
{0.25em}% Space after theorem head
{\small\sffamily\color{ocre}\thmname{#1}\nobreakspace\thmnumber{\@ifnotempty{#1}{}\@upn{#2}}% Theorem text (e.g. Theorem 2.1)
\thmnote{\nobreakspace\the\thm@notefont\sffamily\bfseries\color{black}---\nobreakspace#3}} % Optional theorem note
\makeatother

% Defines the theorem text style for each type of theorem to one of the three styles above
\newcounter{dummy} 
\newcounter{thm}
\newcounter{correction}
\newcounter{qst}
\theoremstyle{ocrenumbox}
\newtheorem{theoremeT}[dummy]{Théorème}
\newtheorem{exerciseT}{Propriété}
\newtheorem{principeT}{Principe}
\theoremstyle{blacknumex}
\newtheorem{exampleT}{Exemple}
\theoremstyle{blacknumexo}
\newtheorem{exo}[thm]{Exercice}
\newtheorem{corr}[correction]{Correction}
\newtheorem{quest}[qst]{Question}
\theoremstyle{blacknumbox}
\newtheorem{vocabulary}{Vocabulary}[section]
\newtheorem{definitionT}{Définition}
\newtheorem{corollaryT}[dummy]{Corollary}
\theoremstyle{ocrenum}
\newtheorem{proofT}[dummy]{Démonstration}


%----------------------------------------------------------------------------------------
%	DEFINITION OF COLORED BOXES
%----------------------------------------------------------------------------------------

\RequirePackage[framemethod=default]{mdframed} % Required for creating the theorem, definition, exercise and corollary boxes

% Theorem box
\newmdenv[skipabove=7pt,
skipbelow=7pt,
backgroundcolor=black!5,
linecolor=ocre,
innerleftmargin=5pt,
innerrightmargin=5pt,
innertopmargin=10pt,
leftmargin=0cm,
rightmargin=0cm,
innerbottommargin=5pt]{tBox}

%Proposition box	  
\newmdenv[skipabove=7pt,
skipbelow=7pt,
rightline=false,
leftline=true,
topline=false,
bottomline=false,
backgroundcolor=ocre!10,
linecolor=ocre,
innerleftmargin=5pt,
innerrightmargin=5pt,
innertopmargin=10pt,
innerbottommargin=3pt,
leftmargin=0cm,
rightmargin=0cm,
linewidth=4pt]{eBox}	

% Definition box
\newmdenv[skipabove=7pt,
backgroundcolor=ocre!4,
skipbelow=7pt,
rightline=false,
leftline=true,
topline=false,
bottomline=false,
linecolor=ocre,
innerleftmargin=5pt,
innerrightmargin=5pt,
innertopmargin=10pt,
leftmargin=0cm,
rightmargin=0cm,
linewidth=4pt,
innerbottommargin=5pt]{dBox}	

% Corollary box
\newmdenv[skipabove=7pt,
skipbelow=7pt,
rightline=false,
leftline=true,
topline=false,
bottomline=false,
linecolor=gray,
backgroundcolor=black!5,
innerleftmargin=5pt,
innerrightmargin=5pt,
innertopmargin=5pt,
leftmargin=0cm,
rightmargin=0cm,
linewidth=4pt,
innerbottommargin=5pt]{cBox}

\newmdenv[skipabove=7pt,
skipbelow=7pt,
backgroundcolor=black!5,
innerleftmargin=5pt,
topline=false,
bottomline=false,
rightline=false,
leftline=false,
innerrightmargin=5pt,
innertopmargin=5pt,
leftmargin=0cm,
rightmargin=0cm,
innerbottommargin=5pt]{xBox}

% Creates an environment for each type of theorem and assigns it a theorem text style from the "Theorem Styles" section above and a colored box from above
\newenvironment{theorem}{\begin{tBox}\begin{theoremeT}}{\end{theoremeT}\end{tBox}}

\newenvironment{exo2}{\noindent \begin{exo}\item\relax \noindent \begin{eBox}\item\relax}{\end{eBox}\end{exo}}


\newenvironment{proposition}{\begin{eBox}\begin{exerciseT}}{\hfill{\color{ocre}}\end{exerciseT}\end{eBox}}		

\newenvironment{principe}{\begin{eBox}\begin{principeT}}{\hfill{\color{ocre}}\end{principeT}\end{eBox}}	
		  
\newenvironment{definition}{\begin{dBox}\begin{definitionT}}{\end{definitionT}\end{dBox}}	

\newenvironment{example}{\begin{xBox}\begin{exampleT}}{\hfill{\tiny\ensuremath{\blacksquare}}\end{exampleT}\end{xBox}}

\newenvironment{demonstration}{\begin{proofT}}{\hfill{\tiny\ensuremath{\square}}\end{proofT}}		
\newenvironment{corollary}{\begin{cBox}\begin{corollaryT}}{\end{corollaryT}\end{cBox}}	

%----------------------------------------------------------------------------------------
%	REMARK ENVIRONMENT
%----------------------------------------------------------------------------------------

\newenvironment{remark}{\par\vspace{5pt}\small % Vertical white space above the remark and smaller font size
\begin{list}{}{
\leftmargin=25pt % Indentation on the left
\rightmargin=15pt}\item\ignorespaces % Indentation on the right
\makebox[-2.5pt]{\begin{tikzpicture}[overlay]
\node[draw=ocre!60,line width=1pt,circle,fill=ocre!25,font=\sffamily\bfseries,inner sep=2pt,outer sep=0pt] at (-15pt,0pt){\textcolor{ocre}{R}};\end{tikzpicture}} % Orange R in a circle
\advance\baselineskip -1pt}{\end{list}\vskip5pt} % Tighter line spacing and white space after remark

%----------------------------------------------------------------------------------------
%	SECTION NUMBERING IN THE MARGIN
%----------------------------------------------------------------------------------------

\makeatletter
\renewcommand{\@seccntformat}[1]{\llap{\textcolor{ocre}{\csname the#1\endcsname}\hspace{1em}}}                    
\renewcommand{\section}{\@startsection{section}{1}{\z@}
{-4ex \@plus -1ex \@minus -.4ex}
{1ex \@plus.2ex }
{\normalfont\LARGE\sffamily\bfseries}}
\renewcommand{\subsection}{\@startsection {subsection}{2}{\z@}
{-3ex \@plus -0.1ex \@minus -.4ex}
{0.5ex \@plus.2ex }
{\normalfont\sffamily\bfseries}}
\renewcommand{\subsubsection}{\@startsection {subsubsection}{3}{\z@}
{-2ex \@plus -0.1ex \@minus -.2ex}
{.2ex \@plus.2ex }
{\normalfont\small\sffamily\bfseries}}                        
\renewcommand\paragraph{\@startsection{paragraph}{4}{\z@}
{-2ex \@plus-.2ex \@minus .2ex}
{.1ex}
{\normalfont\small\sffamily\bfseries}}

%----------------------------------------------------------------------------------------
%	PART HEADINGS
%----------------------------------------------------------------------------------------

% Numbered part in the table of contents
\newcommand{\@mypartnumtocformat}[2]{%
	\setlength\fboxsep{0pt}%
	\noindent\colorbox{ocre!20}{\strut\parbox[c][.7cm]{\ecart}{\color{ocre!70}\Large\sffamily\bfseries\centering#1}}\hskip\esp\colorbox{ocre!40}{\strut\parbox[c][.7cm]{\linewidth-\ecart-\esp}{\Large\sffamily\centering#2}}%
}

% Unnumbered part in the table of contents
\newcommand{\@myparttocformat}[1]{%
	\setlength\fboxsep{0pt}%
	\noindent\colorbox{ocre!40}{\strut\parbox[c][.7cm]{\linewidth}{\Large\sffamily\centering#1}}%
}

\newlength\esp
\setlength\esp{4pt}
\newlength\ecart
\setlength\ecart{1.2cm-\esp}
\newcommand{\thepartimage}{}%
\newcommand{\partimage}[1]{\renewcommand{\thepartimage}{#1}}%
\def\@part[#1]#2{%
\ifnum \c@secnumdepth >-2\relax%
\refstepcounter{part}%
\addcontentsline{toc}{part}{\texorpdfstring{\protect\@mypartnumtocformat{\thepart}{#1}}{\partname~\thepart\ ---\ #1}}
\else%
\addcontentsline{toc}{part}{\texorpdfstring{\protect\@myparttocformat{#1}}{#1}}%
\fi%
\startcontents%
\markboth{}{}%
{\thispagestyle{empty}%
\begin{tikzpicture}[remember picture,overlay]%
\node at (current page.north west){\begin{tikzpicture}[remember picture,overlay]%	
\fill[ocre!20](0cm,0cm) rectangle (\paperwidth,-\paperheight);
\node[anchor=north] at (4cm,-3.25cm){\color{ocre!40}\fontsize{220}{100}\sffamily\bfseries\thepart}; 
\node[anchor=south east] at (\paperwidth-1cm,-\paperheight+1cm){\parbox[t][][t]{8.5cm}{
\printcontents{l}{0}{\setcounter{tocdepth}{1}}% The depth to which the Part mini table of contents displays headings; 0 for chapters only, 1 for chapters and sections and 2 for chapters, sections and subsections
}};
\node[anchor=north east] at (\paperwidth-1.5cm,-3.25cm){\parbox[t][][t]{15cm}{\strut\raggedleft\color{white}\fontsize{30}{30}\sffamily\bfseries#2}};
\end{tikzpicture}};
\end{tikzpicture}}%
\@endpart}
\def\@spart#1{%
\startcontents%
\phantomsection
{\thispagestyle{empty}%
\begin{tikzpicture}[remember picture,overlay]%
\node at (current page.north west){\begin{tikzpicture}[remember picture,overlay]%	
\fill[ocre!20](0cm,0cm) rectangle (\paperwidth,-\paperheight);
\node[anchor=north east] at (\paperwidth-1.5cm,-3.25cm){\parbox[t][][t]{15cm}{\strut\raggedleft\color{white}\fontsize{30}{30}\sffamily\bfseries#1}};
\end{tikzpicture}};
\end{tikzpicture}}
\addcontentsline{toc}{part}{\texorpdfstring{%
\setlength\fboxsep{0pt}%
\noindent\protect\colorbox{ocre!40}{\strut\protect\parbox[c][.7cm]{\linewidth}{\Large\sffamily\protect\centering #1\quad\mbox{}}}}{#1}}%
\@endpart}
\def\@endpart{\vfil\newpage
\if@twoside
\if@openright
\null
\thispagestyle{empty}%
\newpage
\fi
\fi
\if@tempswa
\twocolumn
\fi}

%----------------------------------------------------------------------------------------
%	CHAPTER HEADINGS
%----------------------------------------------------------------------------------------

% A switch to conditionally include a picture, implemented by Christian Hupfer
\newif\ifusechapterimage
\usechapterimagetrue
\newcommand{\thechapterimage}{}%
\newcommand{\chapterimage}[1]{\ifusechapterimage\renewcommand{\thechapterimage}{#1}\fi}%
\newcommand{\autodot}{.}
\def\@makechapterhead#1{%
{\parindent \z@ \raggedright \normalfont
\ifnum \c@secnumdepth >\m@ne
\if@mainmatter
\begin{tikzpicture}[remember picture,overlay]
\node at (current page.north west)
{\begin{tikzpicture}[remember picture,overlay]
\node[anchor=north west,inner sep=0pt] at (0,0) {\ifusechapterimage\includegraphics[width=\paperwidth]{\thechapterimage}\fi};
\draw[anchor=west] (\Gm@lmargin,-3cm) node [line width=2pt,rounded corners=15pt,draw=ocre,fill=white,fill opacity=0.5,inner sep=15pt]{\strut\makebox[22cm]{}};
\draw[anchor=west] (\Gm@lmargin+.3cm,-3cm) node {\huge\sffamily\bfseries\color{black}\thechapter\autodot~#1\strut};
\end{tikzpicture}};
\end{tikzpicture}
\else
\begin{tikzpicture}[remember picture,overlay]
\node at (current page.north west)
{\begin{tikzpicture}[remember picture,overlay]
\node[anchor=north west,inner sep=0pt] at (0,0) {\ifusechapterimage\includegraphics[width=\paperwidth]{\thechapterimage}\fi};
\draw[anchor=west] (\Gm@lmargin,-3cm) node [line width=2pt,rounded corners=15pt,draw=ocre,fill=white,fill opacity=0.5,inner sep=15pt]{\strut\makebox[22cm]{}};
\draw[anchor=west] (\Gm@lmargin+.3cm,-3cm) node {\huge\sffamily\bfseries\color{black}#1\strut};
\end{tikzpicture}};
\end{tikzpicture}
\fi\fi\par\vspace*{50\p@}}}

%-------------------------------------------

\def\@makeschapterhead#1{%
\begin{tikzpicture}[remember picture,overlay]
\node at (current page.north west)
{\begin{tikzpicture}[remember picture,overlay]
\node[anchor=north west,inner sep=0pt] at (0,0) {\ifusechapterimage\includegraphics[width=\paperwidth]{\thechapterimage}\fi};
\draw[anchor=west] (\Gm@lmargin,-3cm) node [line width=2pt,rounded corners=15pt,draw=ocre,fill=white,fill opacity=0.5,inner sep=15pt]{\strut\makebox[22cm]{}};
\draw[anchor=west] (\Gm@lmargin+.3cm,-3cm) node {\huge\sffamily\bfseries\color{black}#1\strut};
\end{tikzpicture}};
\end{tikzpicture}
\par\vspace*{50\p@}}
\makeatother

%----------------------------------------------------------------------------------------
%	LINKS
%----------------------------------------------------------------------------------------

\usepackage{hyperref}
\hypersetup{hidelinks,backref=true,pagebackref=true,hyperindex=true,colorlinks=false,breaklinks=true,urlcolor=ocre,bookmarks=true,bookmarksopen=false}

\usepackage{bookmark}
\bookmarksetup{
open,
numbered,
addtohook={%
\ifnum\bookmarkget{level}=0 % chapter
\bookmarksetup{bold}%
\fi
\ifnum\bookmarkget{level}=-1 % part
\bookmarksetup{color=ocre,bold}%
\fi
}
}

\renewcommand*\thesection{\arabic{section}}

\newcommand*{\coord}[3]{% 
  \ensuremath{\overrightarrow{#1}\, 
    \begin{pmatrix} 
      #2\\ 
      #3 
    \end{pmatrix}}}
    
  \newcommand*{\coordb}[2]{% 
  \ensuremath{ 
    \begin{pmatrix} 
      #1\\ 
      #2 
    \end{pmatrix}}}

\newcommand*{\coorde}[4]{% 
  \renewcommand{\arraystretch}{1}\ensuremath{\overrightarrow{#1}\, 
    \begin{pmatrix} 
      #2\\ 
      #3 \\
      #4
    \end{pmatrix}}}    
  \newcommand*{\coordbe}[3]{% 
 \renewcommand{\arraystretch}{1} \ensuremath{ 
    \begin{pmatrix} 
      #1\\ 
      #2 \\
      #3
    \end{pmatrix}}}  
    
\newcommand{\Card}{\mathrm{Card}}


\DeclareExerciseCollection[topic=lgn01]{lgn01}
\DeclareExerciseCollection{lgn02}
\DeclareExerciseCollection{lgn03}
\DeclareExerciseCollection{lgn04}

\begin{document}

\chapterimage{../../Pictures/background}

\chapter{Cours : Loi des grands nombres}


\section{Opérations sur les variables aléatoires}

\subsection{Sommes et produits par un réel}

\begin{definition}Soit $X$ une variable aléatoires réelle, définie sur $\Omega$. Soit $a$ un réel non nul et $b$ un réel. 

La variable aléatoire $aX+b$ est définie par : pour tout $\omega \in \Omega$, $(aX+b)(\omega) = a \times X(\omega) +b$

Ainsi, pour tout réel $k$, on a $\mathbb{P}(aX+b=k) = \mathbb{P}\left( X = \dfrac{k-b}{a} \right)$.\end{definition}

\begin{example}On considère une variable aléatoire $X$ dont la loi est résumée dans le tableau suivant.

\renewcommand{\arraystretch}{2.2}
\begin{center}
\begin{tabular}{|l|c|c|c|}
\hline
$k$ & $-2$& $3$ & $7$ \\
\hline
$\mathbb{P}(X=k)$ & $\dfrac{1}{2}$ & $\dfrac{1}{3}$ & $\dfrac{1}{6}$\\
\hline \end{tabular}
\end{center}

\begin{minipage}{0.68\linewidth}On note $Y$ la variable aléatoire telle que $Y=3X-1$.\\ 
Puisque $X$ prend les valeurs $-2$, $3$ et $7$, $Y$ prend les valeurs $-7$, $8$ et $20$. 

Par ailleurs, $\mathbb{P}(Y=-7)=\mathbb{P}(3X-1=-7)=\mathbb{P}(X=-2)$.\\
La loi de $Y$ peut être résumée par le tableau ci-contre.\end{minipage}\hfill\begin{minipage}{0.27\linewidth}

\renewcommand{\arraystretch}{2.2}
\begin{tabular}{|l|c|c|c|}
\hline
$k$ & $-7$& $8$ & $20$ \\
\hline
$\mathbb{P}(Y=k)$ & $\dfrac{1}{2}$ & $\dfrac{1}{3}$ & $\dfrac{1}{6}$\\
\hline \end{tabular}
\end{minipage}
\end{example}


\begin{example}On considère le jeu suivant : on paye 10 euros puis on lance 4 dés équilibrés à 6 faces, numérotés de 1 à 6. On remporte alors 6 euros par dé qui tombe sur le nombre 6.

Notons $X$ le nombre de 6 obtenus et $Y$ le gain en euros à l'issue de ce jeu. $X$ suit une loi binomiale de paramètres 4 et $\frac{1}{6}$. 

Par ailleurs, on a $Y=6X-10$. En effet, $-10$ représente le coût fixe du jeu, et $6X$ le gain, 6 euros par numéro 6 obtenu. Il est donc facile d'obtenir la loi de $Y$ à partir de celle de $X$.

\begin{minipage}{0.49\linewidth}
\begin{center}
\renewcommand{\arraystretch}{2.2}
\begin{tabular}{|l|c|c|c|c|c|}
\hline
$k$ & $0$& $1$ & $2$  & $3$ & $4$\\
\hline
$\mathbb{P}(X=k)$ & $\dfrac{625}{1296}$ & $\dfrac{125}{324}$ & $\dfrac{25}{216}$ & $\dfrac{5}{324}$ & $\dfrac{1}{1296}$ \\
\hline \end{tabular}
\end{center}

\end{minipage}\hfill\begin{minipage}{0.49\linewidth}
\begin{center}
\renewcommand{\arraystretch}{2.2}
\begin{tabular}{|l|c|c|c|c|c|}
\hline
$k$ & $-10$& $-4$ & $2$  & $8$ & $14$\\
\hline
$\mathbb{P}(Y=k)$ & $\dfrac{625}{1296}$ & $\dfrac{125}{324}$ & $\dfrac{25}{216}$ & $\dfrac{5}{324}$ & $\dfrac{1}{1296}$ \\
\hline \end{tabular}
\end{center}

\end{minipage}
\vspace{-0,1cm}
\end{example}


\subsection{Variables aléatoires indépendantes}

\begin{definition}Soit $n$ un entier naturel et $X_1$, $X_2$, ..., $X_n$ des variables aléatoires définies sur un univers $\Omega$.

On dit que les variables aléatoires $X_1$, $X_2$, ... $X_n$ sont mutuellement indépendantes (ou tout simplement indépendantes) si, pour tous réels $x_1$, $x_2$, ..., $x_n$, on a
\[\pp(X=x_1 \cap X_2=x_2 \cap \ldots \cap X_n=x_n)=\pp(X_1=x_1) \times \pp(X_2=x_2) \times \ldots \times \pp(X_n=x_n).\]\end{definition}

\begin{example}On lance trois fois de suite un dé équilibré à six faces, numérotées de 1 à 6.On appelle $X$ le numéro obtenu au premier lancer, $Y$ le numéro obtenu au deuxième lancer et $Z$ le numéro obtenu au troisième lancer. Alors, les variables aléatoires $X$, $Y$, et $Z$ sont indépendantes.

Plus généralement, si l'on considère une succession d'épreuves aléatoires indépendantes, chacune étant reliée à une variable aléatoire réelle, alors ces variables aléatoires sont indépendantes.\end{example}



\subsection{Somme de deux variables aléatoires}


\begin{definition}Soit $X$ et $Y$ deux variables aléatoires réelles définies sur $\Omega$. 

La variable aléatoire $Z=X+Y$ est définie par : pour tout $\omega \in \Omega$,  $Z(\omega) =  X(\omega) + Y(\omega)$.

Par ailleurs, pour tout réel $k$, on a $\pp(Z=k) = \displaystyle \sum_{i+j=k} \mathbb{P}((X=i) \cap (Y=j)).$

Si les variables aléatoires $X$ et $Y$ sont indépendantes, on a alors $\mathbb{P}(Z=k)=\displaystyle\sum_{i+j=k} \mathbb{P}(X=i) \mathbb{P}(Y=j)$.\end{definition}

\begin{example} 
Un supporter de football a étudié le nombre de buts marqués par son équipe au cours d'une saison. On considère un match au hasard de cette équipe et on appelle $X$ le nombre de buts marqué par cette équipe en première mi-temps de ce match et $Y$ le nombre de buts marqués en second mi-temps. Ainsi, $X+Y$ représente le nombre de buts marqués par l'équipe en question au cours du match.

D'après l'étude de ce supporter, on peut construire l'arbre de probabilités suivant.

\tikzstyle{level 1}=[level distance=3cm, sibling distance=5.5cm]
\tikzstyle{level 2}=[level distance=4cm, sibling distance=2cm]
\tikzstyle{level 3}=[level distance=4cm, sibling distance=0.3cm]
\tikzstyle{level 4}=[level distance=4cm, sibling distance=0.3cm]

% Define styles for bags and leafs
\tikzstyle{bag} = [text centered]
\tikzstyle{end} = [circle, minimum width=3pt,fill, inner sep=0pt]


\begin{center}
\begin{tikzpicture}[scale=0.7,grow=right,sloped]
\draw (3,8,5) node {\textbf{$X$}};
\draw (7,8,5) node {\textbf{$Y$}};
\draw (15,8,5) node {\textbf{$X+Y$}};
\draw (11,8,5) node {\textbf{$(X,Y)$}};
\node[bag] { }
    child {
        node[bag] {2} 
        child {
                node[bag] {2}
                child {
                node[bag] {(2;2)}
                child {
                node[bag] {$4$}}
                edge from parent[dashed,->,>=latex] node[below] {$ $}
                }
                edge from parent node[above] {$0,2$} 
               }
		child {
                node[bag] {1}
                child {
                node[bag] {(2;1)}
                child {
                node[bag] {$3$}}
                edge from parent[dashed,->,>=latex] node[below] {$ $}
                }
                edge from parent node[above] {$0,5$} 
                }
		child {
                node[bag] {0}
                child {
                node[bag] {(2;0)}
                child {
                node[bag] {$2$}}
                edge from parent[dashed,->,>=latex] node[below] {$ $}
                }
                edge from parent node[above] {$0,3$} 
                }
            edge from parent node[above] {$0,05$} 
    }
    child {
        node[bag] {1} 
        child {
                node[bag] {2}
                child {
                node[bag] {(1;2)}
                child {
                node[bag] {$3$}}
                edge from parent[dashed,->,>=latex] node[below] {$ $}
                }
                edge from parent node[above] {$0,15$} 
               }
		child {
                node[bag] {1}
                child {
                node[bag] {(1;1)}
                child {
                node[bag] {$2$}}
                edge from parent[dashed,->,>=latex] node[below] {$ $}
                }
                edge from parent node[above] {$0,5$} 
                }
		child {
                node[bag] {0}
                child {
                node[bag] {(1;0)}
                child {
                node[bag] {$1$}}
                edge from parent[dashed,->,>=latex] node[below] {$ $}
                }
                edge from parent node[above] {$0,35$} 
                }
            edge from parent node[above] {$0,4$} 
    }
    child {
        node[bag] {0} 
        child {
                node[bag] {2}
                child {
                node[bag] {(0;2)}
                child {
                node[bag] {$2$}}
                edge from parent[dashed,->,>=latex] node[below] {$ $}
                }
                edge from parent node[above] {$0,1$} 
               }
		child {
                node[bag] {1}
                child {
                node[bag] {(0;1)}
                child {
                node[bag] {$1$}}
                edge from parent[dashed,->,>=latex] node[below] {$ $}
                }
                edge from parent node[above] {$0,6$} 
                }
		child {
                node[bag] {0}
                child {
                node[bag] {(0;0)}
                child {
                node[bag] {$0$}}
                edge from parent[dashed,->,>=latex] node[below] {$ $}
                }
                edge from parent node[above] {$0,3$} 
                }
            edge from parent node[above] {$0,55$} 
    }
    ;
\end{tikzpicture}
\end{center}

La variable aléatoire $X+Y$ prend alors les valeurs 0, 1, 2, 3 et 4.

Il est alors possible d'établir la loi de  en s'appuyant sur cet arbre de probabilités. Par exemple, on a
\[\pp(X+Y=1)=\pp(X=0 \cap Y=1)+\pp(X=1 \cap Y=0)=0,55 \times 0,6 + 0,4 \times 0,35=0,47\]

La loi de $X+Y$ peut alors être résumée dans le tableau suivant.

\renewcommand{\arraystretch}{1.5}
\begin{center}
\begin{tabularx}{0.7\linewidth}{|l|X|X|X|X|X|}
\hline
$k$ & 0 & 1 & 2 & 3 & 4 \\
\hline
$\pp(X+Y=k)$ & 0,165 & 0,47 & 0,27 & 0,085&0,01 \\
\hline
\end{tabularx}

\end{center}
\vspace{-0,5cm}
\end{example}



\section{Espérance et variance d'une somme de variables}

\subsection{Cas général}

\begin{proposition}Soit $X$ et $Y$ deux variables aléatoires, $a$ et $b$ deux réels. On a
\[E(aX+b)=aE(X) + b \qquad \text{et} \qquad E(X+Y)=E(X)+ E(Y).\]
\vspace{-0,5cm}\end{proposition}

\begin{example} On considère le jeu suivant : la participation est fixée à 8 euros. On lance un dé équilibré à 6 faces, numérotées de 1 à 6 et on remporte deux fois la somme inscrite sur le dé. On considère la variable aléatoire $X$ qui donne le résultat du lancer et $Y$ le gain du joueur, positif ou négatif. On a alors $Y=2X-8$.

L'espérance de $X$ est 3,5. Ainsi, $\mathbb{E}(Y)= \mathbb{E}(2X-8)=2\mathbb{E}(X)-8=2 \times 3,5 - 8 = -1$.\\ L'espérance étant négative, le jeu est défavorable au joueur.\end{example}

\begin{example}On reprend l'exemple précédent de l'étude du supporter. La loi de $X$ s'obtient simplement à l'aide du premier sous-arbre. La loi de $X+Y$ avait été déterminée.

\begin{minipage}{0.35\linewidth}

\renewcommand{\arraystretch}{1.5}
\begin{tabularx}{\linewidth}{|l|X|X|X|}
\hline
$k$ & 0 & 1 & 2  \\
\hline
$\pp(X=k)$ & 0,55 & 0,4 & 0,05  \\
\hline
\end{tabularx}
\end{minipage}\hfill\begin{minipage}{0.6\linewidth}
\renewcommand{\arraystretch}{1.5}
\begin{tabularx}{\linewidth}{|l|X|X|X|X|X|}
\hline
$k$ & 0 & 1 & 2 & 3 & 4 \\
\hline
$\pp(X+Y=k)$ & 0,165 & 0,47 & 0,27 & 0,085&0,01 \\
\hline
\end{tabularx}

\end{minipage}

On a alors $E[X]= 1 \times 0,4 + 2 \times 0,05 = 0,5$ et $E[X+Y]= 1 \times 0,47 + 2 \times 0,27+3 \times 0,085+4 \times 0,01=1,305$.\\
Or, $E[X+Y]=E[X]+E[Y]$. Ainsi, $E[Y]=E[X+Y]-E[X]=1,305-0,5=0,805$.

Il est aussi possible de déterminer la loi de la variable aléatoire $Y$. En effet, les événements $\{X=0\}$, $\{X=1\}$ et $\{X=2\}$ forment une partition de l'univers. Ainsi, d'après la formule des probabilités totales,
\[\pp(Y=0)=\pp(Y=0 \cap X=0)+\pp(Y=0 \cap X=1)+\pp(Y=0\cap X=2)=0,55 \times 0,3 + 0,4 \times 0,35 + 0,05 \times 0,3.\]
Ainsi, $\pp(Y=0)=0,32$. On calcule de la même manière $\pp(Y=1)$ et $\pp(Y=2)$ pour obtenir le tableau suivant.
\begin{center}
\renewcommand{\arraystretch}{1.5}
\begin{tabularx}{0.6\linewidth}{|l|X|X|X|}
\hline
$k$ & 0 & 1 & 2  \\
\hline
$\pp(Y=k)$ & 0,32 & 0,555 & 0,125  \\
\hline
\end{tabularx}
\end{center}

On retrouve alors que $E[Y] = 0 \times 0,32 + 1 \times 0,555 + 2 \times 0,125 = 0,805$.
\end{example}

\textbf{Remarque} : Dans l'exemple précédent, les variables aléatoires $X$ et $Y$ n'étaient pas indépendantes, et pourtant, l'espérance de la somme pouvait être exprimée comme la somme des espérances.

L'hypothèse d'indépendance des variables aléatoires est toutefois primordiales pour la propriété qui suit.

\begin{proposition}Soit $X$ et $Y$ deux variables aléatoires \textbf{indépendantes}, $a$ un réel.
\[V(aX)=a^2\times V(X) \qquad\text{et}\qquad  V(X+Y)=V(X)+V(Y).\]\vspace{-0,5cm}
\end{proposition}

Il est important de ne pas oublier le carré : une variance est toujours positive ! 

Attention également à ne pas se faire piéger lorsque l'on calcule la variance de la somme. \\
Si $X$ et $Y$ sont deux variables aléatoire réelles indépendantes, alors \[V(X-Y)=V(X)+V(-Y)=V(X)+(-1)^2V(Y)=V(X)+V(Y).\]

\begin{example}Toujours dans l'exemple précédente, on peut comparer les variances de $X$, $Y$ et $X+Y$. 

Calculons la variance de $X$. On rappelle pour cela la formule de Koenig-Huygens : $V(X)=E[X^2]-(E[X])^2$.

\renewcommand{\arraystretch}{1.5}
\begin{minipage}{0.45\linewidth}
\begin{tabularx}{\linewidth}{|l|X|X|X|}
\hline
$k$ & 0 & 1 & 2  \\
\hline
$\pp(X=k)$ & 0,55 & 0,4 & 0,05  \\
\hline
\end{tabularx}
\end{minipage}\hfill\begin{minipage}{0.45\linewidth}
\begin{tabularx}{\linewidth}{|l|X|X|X|}
\hline
$k$ & 0 & 1 & 4  \\
\hline
$\pp(X^2=k)$ & 0,55 & 0,4 & 0,05  \\
\hline
\end{tabularx}
\end{minipage}

Ainsi, $E[X^2]=0 \times 0,55 + 1 \times 0,4 + 4 \times 0,05 = 0,6$ et donc $V(X)=0,6-0,5^2=0,35$.

De même, $E[Y^2]=0 \times 0,32 + 1 \times 0,555 + 4 \times 0,125 = 1,055$ et donc $V(Y)=1,055-0,805^2 = 0,406975$.

Enfin, $E[(X+Y)^2]=0\times 0,165 + 1 \times 0,47 + 4 \times 0,27 + 9 \times 0,085 + 16 \times 0,01 = 2,475$\\
On obtient donc $V(X+Y)=2,475-1,305^2=0,771975$.

En particulier, on voit que $V(X+Y)\neq V(X) + V(Y)$. Ce n'est pas surprenant puisque les variables aléatoires $X$ et $Y$ ne sont pas indépendantes.

\end{example}

\subsection{Applications à la loi binomiale}

\begin{proposition} Soit $X_1$, $X_2$, ..., $X_n$ des variables aléatoires de Bernoulli indépendantes de paramètre $p$. 

Notons $S=X_1+X_2+\ldots + X_n$. La variable aléatoire $S$ suit une loi binomiale de paramètres $n$ et $p$.\end{proposition}


\begin{proposition}Soit $X$ une variable aléatoire qui suit une loi binomiale $\mathcal{B}(n,p)$. L'espérance, la variance et l'écart-type de $X$ valent respectivement \[E[X]=np, \quad V(X)=np(1-p), \quad \sigma(X)=\sqrt{np(1-p)}\]\end{proposition}

\begin{demonstration}Soit $X$ une variable aléatoire qui suit une loi binomiale $\mathcal{B}(n,p)$. 

On considère des variables aléatoires indépendantes $X_1$, $X_2$, ..., $X_n$  suivant toute une loi de Bernoulli de paramètre $p$. Pour tout entier naturel $i$ compris entre 1 et $p$, on a en particulier $E[X_i]=p$ et $V(X_i)=p(1-p)$.

On note alors $S=X_1+X_2+\ldots +X_n$. D'après la propriété précédente, $S$ suit également une loi binomiale $\mathcal{B}(n,p)$. $S$ et $X$ sont donc de même loi et ont la même espérance.

Or, on a $E[S] = p + p + \ldots + p = np$.\\ De plus, les variables $X_i$ étant indépendants, $V(S)=p(1-p)+p(1-p)+\ldots + p(1-p)=np(1-p)$.\end{demonstration}

\begin{example}On lance 12 dés équilibrées à 6 faces, numérotées de 1 à 6. On note $Y$ la variable aléatoires qui compte le nombre de 6 obtenus. 

$Y$ suit une loi binomiale de paramètres $12$ et $\dfrac{1}{6}$. Ainsi, $E[Y]=12 \times \dfrac{1}{6}=2$ et $V(Y)=12 \times \dfrac{1}{6}\times \dfrac{5}{6}=\dfrac{5}{3}$\end{example}



\section{Concentration et loi des grand nombres}
\subsection{Échantillon de variables aléatoires}

\begin{definition} Un échantillon est un ensemble de variables aléatoires réelles $(X_1, ..., X_n)$ indépendantes et de même loi.

La variable aléatoire moyenne de cette échantillon est la variable aléatoire notée $M_n$ ou $\overline{X}$, définie par 
\[ M_n = \dfrac{1}{n} (X_1+X_2+\ldots + X_n)\]\end{definition}


\begin{proposition}On a alors $\mathbb{E}(M_n)=E[X_1]$, $V(M_n)=\dfrac{1}{n}V(X_1)$ et $\sigma(M_n)=\dfrac{1}{\sqrt{n}}\sigma(X_1)$\end{proposition}

\begin{demonstration}On a en effet
\[ E[M_n]=E\left[\dfrac{1}{n} (X_1+X_2+\ldots + X_n)\right]=\dfrac{1}{n}(E[X_1]+E[X_2]+\dots + E[X_n])=\dfrac{1}{n} \times nE[X_1] = E[X_1]\]
Par ailleurs, les variables aléatoires étant indépendantes, 
\[ V(M_n)=V\left(\dfrac{1}{n} (X_1+X_2+\ldots + X_n)\right)=\dfrac{1}{n^2}(V(X_1)+V(X_2)+\dots + V(X_n))\]
et donc 
\[V(M_n)=\dfrac{1}{n^2} \times nV(X_1) = \dfrac{V(X_1)}{n}\]
\end{demonstration}


\begin{example} On considère une variable aléatoire $X$ qui suit une loi binomiale de paramètre $3$ et $\dfrac{1}{3}$. 

On rappelle que $E[X]=3\times \dfrac{1}{3}=1$ et $V(X)=3\times \dfrac{1}{3} \times \left(1-\dfrac{1}{3}\right)=\dfrac{2}{3}$.

On considère un échantillon $(X_1, \ldots, X_{100})$ de variables aléatoires indépendantes de même loi que $X$ et on note $M_{100} = \dfrac{1}{100} (X_1+X_2+\ldots + X_{100})$.

On a alors $E[M_{100}]=E[X]=1$ et $V(M_n)=\dfrac{V(X)}{100}=\dfrac{2}{300}$.\end{example}

On remarque que lorsque $n$ tend vers $+\infty$, la variance de $M_n$ tend vers 0 alors que l'espérance ne change pas. Cela signifie intuitivement que la variable aléatoire $M_n$ se rapproche d'une variable aléatoire "constante". Ce comportement sera précisé plus en détails dans les parties qui suivent.

\newpage 
\subsection{Inégalité de Bienaymé-Tchebychev}

\begin{proposition}[Inégalité de Bienaymé-Tchebychev] Soit $X$ une variable aléatoire réelle. Pour tout réel $\delta$ strictement positif
\[ \mathbb{P}(|X-E[X]| \geqslant \delta) \leqslant \dfrac{V(x)}{\delta ^2}\]
\vspace{-0,5cm}\end{proposition}

Cette inégalité illustre le fait que la variance permet de mesurer l'écart d'une variable aléatoire par rapport à son espérance.

\begin{example}Soit $X$ une variable aléatoire d'espérance 10 et de variance 1.

D'après l'inégalité de Bienaymé-Tchebychev, appliquée à $\delta = 4$, on a $\pp(|X-10|\geqslant 4) \leqslant \dfrac{1}{4^2}$.

Par ailleurs, les événements $|X-10|\geqslant 4$ et $|X-10|<4$ étant contraires, on a donc \[\pp(|X-10|< 4) = 1 - \pp(|X-10|\geqslant 4) \geqslant 1-\dfrac{1}{4^2}.\]

Or, $|X-10|<4$ est équivalent à $X \in ]10-4 ; 10+4[$, c'est-à-dire $X\in ]6 ; 14[$. Finalement, $\pp(X\in ]6 ;14[) \geqslant \dfrac{15}{16}$.\end{example}

\begin{example} On lance 180 fois un dé équilibré à 6 faces, numérotées de 1 à 6. On appelle $X$ la variable aléatoire qui donne le nombre de 1 obtenus. $X$ suit donc une loi binomiale de paramètres $180$ et $\dfrac{1}{6}$.

Ainsi, $E(X)=180 \times \dfrac{1}{6}=30$ et $V(X)=180 \times \dfrac{1}{6} \times \dfrac{5}{6}=25$. 

Lors des précédentes chapitres, nous avons interprété l'espérance comme une moyenne si l'on répète un grand nombre de fois une expérience aléatoire. Ainsi, si on lance 180 fois un dé, on s'attend en moyenne à avoir 30 fois le nombre 6.

Seulement, tout ceci n'est qu'une moyenne, et il est rare de tomber exactement 30 fois sur la face numéro 6. Cependant, il y a une grande probabilité que le nombre de fois que nous obtenons ce nombre 6 soit proche de 30, et l'inégalité de Bienaymé-Tchebychev peut nous fournir une minoration de cette probabilité.


On souhaite par exemple minorer la probabilité que $X$ soit compris entre 21 et 39. 

Or, $X\in[21 ;39]$ revient à dire que $X \in [30-9 ; 30+9]$ c'est-à-dire $|X-30| \leqslant 9$. Puisque la variable aléatoire $X$ prend des valeurs entières uniquement, ceci est équivalent à $|X-30|<10$.

D'après l'inégalité de Bienaymé-Tchebychev, 
\[ \mathbb{P}( |X-E(X)| \geqslant 10) \leqslant \dfrac{V(X)}{10^2} =\dfrac{25}{100}=\dfrac{1}{4}\]

Ainsi, puisque $\mathbb{P}( |X-E(X)| <10)+\mathbb{P}( |X-E(X)|\geqslant 10)=1$, on a que \[\mathbb{P}( |X-E(X)| < 10)=1-\mathbb{P}( |X-E(X)| \geqslant 10) \geqslant 1- \dfrac{1}{4}\]

Ainsi, $\mathbb{P}( |X-E(X)| < 10) \geqslant \dfrac{3}{4}$.

Si l'on lance 180 dés, la probabilité d'avoir entre 21 et 39 fois le nombre 6 est supérieure à 0,75.\end{example}

Cette borne n'est pas toujours optimale. En l'occurrence, en faisant les calculs de manière exhaustive, on s'aperçoit que $\mathbb{P}( |X-E(X)| < 10) \simeq 0,9434$, mais ce calcul est un poil plus compliqué...



\subsection{Inégalité de concentration}

\begin{proposition}[Inégalité de concentration] Soit $(X_1,\ldots,X_n)$ un échantillon de $n$ variables aléatoires indépendantes, et $M_n$ la variable aléatoire moyenne de cette échantillon. Alors, pour tout réel $\delta$ strictement positif,
\[ \mathbb{P}(|M_n-E(X_1)|\geqslant \delta ) \leqslant \dfrac{V(X_1)}{n\delta^2}\]\end{proposition}

\begin{demonstration} On applique simplement l'inégalité de Bienaymé-Tchebychev à la variable aléatoire $M_n$. Son espérance vaut $E(X_1)$ et sa variance $\dfrac{V(X_1)}{n}$.\end{demonstration}

\begin{example} Soit $X$ une variable aléatoire d'espérance 3 et de variance 100.

 On considère un échantillon $(X_1, \ldots, X_n)$ de variables aléatoires indépendantes de même loi que $X$ et on note $M_n = \dfrac{1}{n} (X_1+X_2+\ldots + X_n)$.

Pour tout entier naturel non nul $n$ et tout réel $\delta$ strictement positif, on a alors $\mathbb{P}(|M_n-E(X_1)|\geqslant \delta ) \leqslant \dfrac{V(X_1)}{n\delta ^2}$,

c'est-à-dire $\mathbb{P}(|M_n-3|\geqslant \delta ) \leqslant \dfrac{100}{n \delta ^2}$.

En particulier, pour $n=100000$ et $\delta=0,1$, on a $\mathbb{P}(|M_n-3|\geqslant 0,1 ) \leqslant \dfrac{100}{100000 \times 0,1^2}$

Ainsi, $\mathbb{P}(|M_n-3|\geqslant 0,1 ) \leqslant 0,1$.

En passant au complémentaire, on obtient alors que $\mathbb{P}(|M_n-3| < 0,1 ) = 1 - \mathbb{P}(|M_n-3|\geqslant 0,1 ) \geqslant 0,9$.

Bien que la variable aléatoire $X$ ait une grande variance, si l'on répète un grand nombre de fois l'expérience aléatoire, la moyenne des résultats est très proche de l'espérance de $X$ : avec probabilité 0,9, la moyenne est entre 2,9 et 3,1.\end{example}


\subsection{Loi des grands nombres}

\begin{theorem}[Loi faible des grands nombres] Soit $(X_1,...,X_n)$ un échantillon de $n$ variables aléatoires indépendantes et $M_n$ la variable aléatoire moyenne de cet échantillon. Pour tout réel $\delta$ strictement positif,
\[ \displaystyle \lim_{n \to + \infty} \mathbb{P}(|M_n-E(X_1)|\geqslant \delta )=0\]\end{theorem}

\begin{demonstration}On applique l'inégalité de concentration à cet échantillon :
\[ \mathbb{P}(|M_n-E(X_1)|\geqslant \delta ) \leqslant \dfrac{V(X_1)}{n\delta^2}\]

Or, $\displaystyle \lim _{n \to + \infty} \dfrac{V(X_1)}{n\delta^2}=0$. De plus, $\mathbb{P}(|M_n-E(X_1)|\geqslant \delta ) \geqslant 0$. 

D'après le théorème d'encadrement, on a donc  $\displaystyle \lim_{n \to + \infty} \mathbb{P}(|M_n-E(X_1)|\geqslant \delta )=0$.\end{demonstration}

\chapter{Exercices}


\section*{Opérations sur les variables aléatoires}

\begin{exercise}On considère la variable aléatoire $X$ dont la loi est résumée dans le tableau suivant.

\begin{minipage}{0.35\linewidth}
\renewcommand{\arraystretch}{2.2}
\begin{center}
\begin{tabular}{|l|c|c|c|c|}
\hline
$k$ & $-3$& $-1$ & $2$ & $4$ \\
\hline
$\mathbb{P}(X=k)$ & $\dfrac{1}{5}$ & $\dfrac{2}{5}$ & $\dfrac{1}{10}$   & $\dots$ \\
\hline \end{tabular}
\end{center}\end{minipage}\begin{minipage}{0.58\linewidth}

\begin{enumerate}
\item Compléter ce tableau avec la probabilité manquante.
\item Donner la loi de la variable aléatoire $Y=X+2$.
\item Donner la loi de la variable aléatoire $Z=2X-1$.
\end{enumerate}\end{minipage}\end{exercise}

\begin{solution}La probabilité manquante est $\dfrac{3}{10}$. On a alors les tableaux suivants pour les lois de $Y$ et $Z$.

\begin{minipage}{0.45\linewidth}
\renewcommand{\arraystretch}{2.2}
\begin{center}
\begin{tabular}{|l|c|c|c|c|}
\hline
$k$ & $-1$& $1$ & $4$ & $6$ \\
\hline
$\mathbb{P}(Y=k)$ & $\dfrac{1}{5}$ & $\dfrac{2}{5}$ & $\dfrac{1}{10}$   & $\dfrac{3}{10}$ \\
\hline \end{tabular}
\end{center}\end{minipage}\hfill \begin{minipage}{0.45\linewidth}
\renewcommand{\arraystretch}{2.2}
\begin{center}
\begin{tabular}{|l|c|c|c|c|}
\hline
$k$ & $-5$& $-1$ & $5$ & $9$ \\
\hline
$\mathbb{P}(Z=k)$ & $\dfrac{1}{5}$ & $\dfrac{2}{5}$ & $\dfrac{1}{10}$   & $\dfrac{3}{10}$ \\
\hline \end{tabular}
\end{center}\end{minipage}

\end{solution}



\begin{exercise}Une entreprise commercialise des lave-vaisselles et propose à ses clients une extension de garantie de 3 ans supplémentaires au prix de 65 euros. Si une panne irréparable survient au cours de cette période, l'entreprise remboursera alors les 399 euros correspondant au prix du lave-vaisselle. D'après les statistiques relevées par l'entreprise, 11,5\% des lave-vaisselles tombent en panne durant cette période de 3 ans. Un client achète un lave-vaisselle avec extension de garantie.
\begin{enumerate}
\item On note $X$ la variable aléatoire qui vaut 1 si le lave-vaisselle tombe en panne durant la période concernant l'extension de garantie et 0 sinon. Quelle est la loi de $X$ ?
\item On note $Y$ le gain réalisée par l'entreprise grâce à l'extension de garantie. Exprimer $Y$ en fonction de $X$.
\end{enumerate}\end{exercise}

\begin{solution}$X$ suit une loi de Bernoulli de paramètre 0,115. Par ailleurs, on a $Y=65-399X$.\end{solution}



\begin{exercise}[subtitle={(Asie 2022)}]Les compagnies aériennes vendent plus de billets qu'il n'y a de places dans les avions car certains passagers ne se présentent pas à l'embarquement du vol sur lequel ils ont réservé. On appelle cette pratique le surbooking. Au vu des statistiques des vols précédents, la compagnie aérienne estime que chaque passager a 5\% de chance de ne pas se présenter à l'embarquement.

Considérons un vol dans un avion de 200 places pour lequel 206 billets ont été vendus. On suppose que la présence à l'embarquement de chaque passager est indépendante des autres passagers et on appelle $X$ la variable aléatoire qui compte le nombre de passagers se présentant à l'embarquement.

\begin{enumerate}
\item Justifier que $X$ suit une loi binomiale et donner ses paramètres.
\item Calculer $\mathbb{P}(X \leqslant 200)$ et interpréter ce résultat dans le contexte de l'exercice.
\item La compagnie aérienne vend chaque billet à 250 euros. Si plus de 200 passagers se présentent à l'embarquement, la compagnie doit rembourser le billet d'avion et payer une pénalité de 600 euros à chaque passager lésé.

On appelle $Y$ la variable aléatoire égale au nombre de passagers qui ne peuvent pas embarquer bien qu'ayant acheté un billet et $C$ la variable aléatoire qui totalise le chiffre d'affaire de la compagnie aérienne sur ce vol. On admet que Y suit la loi de probabilité donnée par le tableau suivant.

\begin{center}
\begin{tabular}{|l|c|c|c|c|c|c|c|}
\hline
$k$ & $0$& $1$ & $2$ & $3$ & $4$ & $5$ & $6$ \\
\hline
$\mathbb{P}(Y=k)$ & $0,94775$ & $0,03063$ & $0,01441$   & $0,00539$ & $0,00151$ & $0,00028$ & \\
\hline \end{tabular}
\end{center}

\begin{enumerate}
\item Compléter la loi de probabilité donnée ci-dessus.
\item Exprimer $C$ en fonction de $Y$ puis donner la loi de la variable aléatoire $C$ sous forme d'un tableau.
\item Calculer l'espérance de $C$ à l'euro près.
\item Comparer le chiffre d'affaires obtenu en vendant exactement 200 billets et le chiffre d'affaires moyen obtenu en pratiquant le surbooking.
\end{enumerate}
\end{enumerate}\end{exercise}

\begin{solution}Les présences des passagers étant indépendantes, $X$ suit une loi binomiale de paramètres 206 et 0,95. 

D'après la calculatrice, $\mathbb{P}(X \leqslant 200) \simeq 0,948$. Il y a donc environ 5,2\% de chances que trop de passagers se présentent à l'embarquement.

On a $\mathbb{P}(Y=6)=1-(\mathbb{P}(Y=0)+\dots + \mathbb{P}(Y=5))=0,00003$.

En vendant 206 billets, la compagnie encaisse $206 \times 250 = 51500$ euros. Seulement, elle doit rembourser 850 euros par client ne pouvant embarquer (représenté par la variable $Y$). Ainsi, $C=51500-850Y$.

\begin{center}
\begin{tabular}{|l|c|c|c|c|c|c|c|}
\hline
$k$ & $51500$& $50650$ & $49800$ & $48950$ & $48100$ & $47250$ & $46400$ \\
\hline
$\mathbb{P}(C=k)$ & $0,94775$ & $0,03063$ & $0,01441$   & $0,00539$ & $0,00151$ & $0,00028$ & \\
\hline \end{tabular}
\end{center}

On a alors $E[C]\simeq 51429$. Si la compagnie avait vendu seulement 200 billets, son chiffre d'affaires aurait été de 50000 euros, le surbooking lui est donc avantageux.\end{solution}





\begin{exercise}On considère deux variables aléatoires $X$ et $Y$ indépendantes dont les lois sont résumées dans les tableaux suivants.
\vskip10pt
\renewcommand{\arraystretch}{2.2}
\hfill
\begin{tabular}{|l|c|c|c|}
\hline
$k$ & $1$& $3$ & $4$  \\
\hline
$\mathbb{P}(X=k)$ & $\dfrac{1}{3}$ & $\dfrac{1}{4}$ & $\ldots$ \\
\hline \end{tabular}\hfill
et
\hfill
\begin{tabular}{|l|c|c|c|}
\hline
$k$ & $1$& $2$ & $3$ \\
\hline
$\mathbb{P}(Y=k)$ & $\dfrac{1}{5}$ & $\dfrac{1}{5}$ & $\ldots$ \\
\hline \end{tabular}\hfill~
\vskip10pt
\begin{enumerate}
\item Compléter ces tableaux avec les probabilités manquantes.
\item Construire le tableau résumant la loi de la variable aléatoire $Z=2X$.
\item Que vaut $\mathbb{P}(X+Y=5)$ ?
\item Construire le tableau résumant la loi de la variable aléatoire $W=X+Y$
\item Construire le tableau résumant la loi de la variable aléatoire $A=3X-2Y$.
\end{enumerate}\end{exercise}

\begin{solution}On complète les tableaux donnant les lois de $X$ et $Y$ en faisant en sorte que la somme des probabilités vaille 1.

\renewcommand{\arraystretch}{2.2}

\begin{center}
 \begin{tabular}{|l|c|c|c|}
\hline
$k$ & $1$& $3$ & $4$  \\
\hline
$\mathbb{P}(X=k)$ & $\dfrac{1}{3}$ & $\dfrac{1}{4}$ & $\dfrac{5}{12}$ \\
\hline \end{tabular}
$\quad$
et
$\quad$
\begin{tabular}{|l|c|c|c|}
\hline
$k$ & $1$& $2$ & $3$ \\
\hline
$\mathbb{P}(Y=k)$ & $\dfrac{1}{5}$ & $\dfrac{1}{5}$ & $\dfrac{3}{5}$ \\
\hline \end{tabular}
\end{center}

Le tableau résumant la loi de la variable aléatoire $Z=2X$ est le suivant.
\begin{center}
\begin{tabular}{|l|c|c|c|}
\hline
$k$ & $2$& $6$ & $8$  \\
\hline
$\mathbb{P}(X=k)$ & $\dfrac{1}{3}$ & $\dfrac{1}{4}$ & $\dfrac{5}{12}$ \\
\hline \end{tabular}\end{center}

L'événement $X+Y=5$ est réalisé lorsque $(X=3 \cap Y=2)$ ou $(X=4 \cap Y=1)$. Ces variables étant indépendantes, on a donc 
\[\mathbb{P}(X+Y=5) = \dfrac{1}{4} \times \dfrac{1}{5} + \dfrac{5}{12} \times \dfrac{1}{5} = \dfrac{4}{30}.\]
Le tableau résumant la loi de la variable aléatoire $W=X+Y$ est le suivant.

\begin{center}
\begin{tabular}{|l|c|c|c|c|c|c|}
\hline
$k$ & $2$& $3$ & $4$ & $5$ & $6$ & $7$\\
\hline
$\mathbb{P}(W=k)$ & $\dfrac{1}{15}$ & $\dfrac{1}{15}$ & $\dfrac{1}{4}$ & $\dfrac{4}{30}$ & $\dfrac{7}{30}$ & $\dfrac{1}{4}$\\
\hline \end{tabular}\end{center}
Le tableau résumant la loi de la variable aléatoire $A=3X-2Y$ est le suivant.
\begin{center}
\begin{tabular}{|l|c|c|c|c|c|c|c|c|c|}
\hline
$k$ & $-3$& $-1$ & $1$ & $3$ & $5$ & $6$ & $7$ & $8$ & $10$ \\
\hline
$\mathbb{P}(A=k)$ & $\dfrac{1}{5}$ & $\dfrac{1}{15}$ & $\dfrac{1}{15}$ & $\dfrac{3}{20}$ & $\dfrac{1}{20}$ & $\dfrac{1}{4}$ & $\dfrac{1}{20}$ & $\dfrac{1}{12}$ & $\dfrac{1}{12}$\\
\hline \end{tabular}\end{center}
\end{solution}




\begin{exercise}On considère deux variables aléatoires $X$ et $Y$ indépendantes dont les lois sont résumées dans les tableaux suivants.
\vskip10pt
\renewcommand{\arraystretch}{2.2}
\hfill
\begin{tabular}{|l|c|c|}
\hline
$k$ & $2$& $3$  \\
\hline
$\mathbb{P}(X=k)$ & $\dfrac{1}{3}$ & $\dfrac{2}{3}$ \\
\hline \end{tabular}
\hfill et \hfill
\begin{tabular}{|l|c|c|c|}
\hline
$k$ & $1$& $4$ & $5$ \\
\hline
$\mathbb{P}(Y=k)$ & $\dfrac{1}{4}$ & $\dfrac{1}{2}$ & $\dfrac{1}{4}$ \\
\hline \end{tabular}\hfill~
\vskip10pt
\begin{enumerate}
\item Construire le tableau résumant la loi de la variable aléatoire $X+Y$.
\item Construire le tableau résumant la loi de la variable aléatoire $2X+3Y$.
\end{enumerate}\end{exercise}

\begin{solution}La loi de $X+Y$ est la suivante.
\begin{center}
\renewcommand{\arraystretch}{2.2}
\begin{tabular}{|l|c|c|c|c|c|}
\hline
$k$ & $3$& $4$ & $6$ & $7$ & $8$  \\
\hline
$\mathbb{P}(X+Y=k)$ & $\dfrac{1}{12}$ & $\dfrac{1}{6}$ & $\dfrac{1}{6}$ & $\dfrac{5}{12}$ & $\dfrac{1}{6}$\\
\hline \end{tabular}\end{center}

Le cas le plus compliqué ici est $X+Y=7$ puisqu'il y a deux cas à étudier. 
\[\mathbb{P}(X+Y=7)=\mathbb{P}(X=2 \cap Y=5)+\mathbb{P}(X=3 \cap Y=4).\]
$X$ et $Y$ étant indépendantes, on a alors
\[\mathbb{P}(X+Y=7)=\mathbb{P}(X=2)\mathbb{P}(Y=5)+\mathbb{P}(X=3)\mathbb{P}( Y=4)=\dfrac{1}{3} \times \dfrac{1}{4}+ \dfrac{2}{3} \times \dfrac{1}{2}=\dfrac{5}{12}.\]

La loi de $2X+3Y$ est la suivante.

\begin{center}
\renewcommand{\arraystretch}{2.2}
\begin{tabular}{|l|c|c|c|c|c|c|}
\hline
$k$ & $7$& $9$ & $16$ & $18$ & $19$ & $21$  \\
\hline
$\mathbb{P}(X+Y=k)$ & $\dfrac{1}{12}$ & $\dfrac{1}{6}$ & $\dfrac{1}{6}$ & $\dfrac{1}{3}$ & $\dfrac{1}{12}$ & $\dfrac{1}{6}$\\
\hline \end{tabular}\end{center}\end{solution}



\section*{Espérance et variance d'une somme de variables}

\begin{exercise}On considère les deux variables aléatoires $X$ et $Y$ de l'exercice précédent.
\begin{enumerate}
\item Calculer l'espérance et la variance de $X$ et de $Y$.
\item En déduire l'espérance et la variance de $3X+2$, de $X+Y$ et de $5X-2Y$.
\end{enumerate}\end{exercise}

\begin{solution}On a $E[X]=2 \times \dfrac{1}{3} + 3 \times \dfrac{2}{3} = \dfrac{8}{3}$ et $E[Y]=1 \times \dfrac{1}{4}+4 \times \dfrac{1}{2} + 5 \times \dfrac{1}{4}= \dfrac{7}{2}$.

Ainsi, $E[3X+2]=3E[X]+2=10$, $E[X+Y]=E[X]+E[Y]=\dfrac{8}{3}+\dfrac{7}{2}= \dfrac{37}{6}$ et \\$E[5X-2Y]=5E[X]-2E[Y]=\dfrac{40}{3}-7=\dfrac{19}{3}$.\end{solution}




\begin{exercise}
Soit $X$ et $Y$ deux variables aléatoire indépendantes telles que $E(X)=3$, $E(Y)=-5$, $V(X)=1$ et $V(Y)=2$.
\begin{enumerate}
\item On considère la variable aléatoire $Z_1=2X+3Y$. Donner l'espérance et la variance de $Z_1$
\item On considère la variable aléatoire $Z_2=4X-2Y$. Donner l'espérance et la variance de $Z_2$.
\item On considère lavariable aléatoire $Z_3=3Y-2X+7$. Donner l'espérance et la variance de $Z_3$.
\end{enumerate}\end{exercise}

\begin{solution}On a $E(Z_1)=2E(X)+3E(Y)=2 \times 3 + 3 \times (-5)=-9$. De plus, les variables $X$ et $Y$ étant indépendantes, $V(Z_1)=V(2X)+V(3Y)=4V(X)+9V(Y)=4 \times 1 + 9 \times 2 = 22$.

On a $E(Z_2)=4E(X)-2E(Y)=4 \times 3 -2 \times (-5)=22$. De plus, les variables $X$ et $Y$ étant indépendantes, $ V(Z_2)=V(4X)+V(-2Y)=16V(X)+4V(Y)=16 \times 1 + 4 \times 2 = 24$.

On a $E(Z_3)=3E(Y)-2E(X)+7=3 \times (-5) -2 \times 3 + 7 = -14$. De plus, les variables $X$ et $Y$ étant indépendantes, et $7$ étant un réel, $V(Z_3)=V(3Y)+V(-2X)+V(7)=9V(Y)+4V(X)+0=9 \times 2 + 4 \times 1 = 22$.\end{solution}




\begin{exercise}On dit qu'une variable aléatoire $X$ est centrée réduite si son espérance est nulle et sa variance vaut 1. 

Montrer que pour toute variable aléatoire $X$ non constante et admettant une espérance et une variance, la variable $Y=\dfrac{X-\mathbb{E(X)}}{\sigma (X)}$ est centrée réduite.\newpage \end{exercise}

\begin{solution}On a $E[Y]=\dfrac{E[X]-E[X]}{\sigma(X)}=0$ et $V(Y) = \dfrac{V(X)+0}{\sigma(X)^2}=\dfrac{V(X)}{V(X)}=1$. $Y$ est centrée et réduite.\end{solution}





\begin{exercise}[subtitle={(Centres étrangers 2024)}]

Un sac opaque contient huit jetons numérotés de 1 à 8, indiscernables au toucher.

À trois reprises, un joueur pioche un jeton dans ce sac, note son numéro, puis le remet dans le sac.

Dans ce contexte, on appelle « tirage » la liste ordonnée des trois numéros obtenus.

Par exemple, si le joueur pioche le jeton numéro 4, puis le jeton numéro 5, puis le jeton numéro 1, alors le tirage correspondant est (4; 5; 1).
\begin{enumerate}
\item Déterminer le nombre de tirages possibles.
\item \begin{enumerate}
\item Déterminer le nombre de tirages sans répétition de numéro.
\item En déduire le nombre de tirages contenant au moins une répétition de numéro.\end{enumerate}\end{enumerate}
On note $X_1$ la variable aléatoire égale au numéro du premier jeton pioché, $X_2$ celle égale au numéro du deuxième jeton pioché et $X_3$ celle égale au numéro du troisième jeton pioché.

Puisqu'il s'agit d'un tirage avec remise, les variables aléatoires $X_1$ ,$X_2$ et $X_3$ sont indépendantes et suivent la même loi de probabilité.

\begin{enumerate}
\setcounter{enumi}{2}
\item Établir la loi de probabilité de la variable aléatoire $X_1$.
\item Déterminer l'espérance de la variable aléatoire $X_1$.\end{enumerate}

On note $S = X_1 + X_2 + X_3$ la variable aléatoire égale à la somme des numéros des trois jetons piochés.

\begin{enumerate}
\setcounter{enumi}{4}
\item Déterminer l'espérance de la variable aléatoire $S$.
\item Déterminer $\pp(S=24)$.
\item Si un joueur obtient une somme supérieure ou égale à 22, alors il gagne un lot.
\begin{enumerate}
\item Justifier qu'il existe exactement 10 tirages permettant de gagner un lot.
\item En déduire la probabilité de gagner un lot.\end{enumerate}\end{enumerate}

\end{exercise}

\begin{solution}\hspace{0pt}

\begin{enumerate}
\item Le nombre de tirages est de $8 \times 8 \times 8=512$.
\item \begin{enumerate}
\item Le nombre de tirages sans répétition est de $8 \times 7 \times 6=336$.
\item Le nombre de tirages avec au moins une répétition est donc $512-336=176$.
\end{enumerate}
\item $X_1$ prend les valeurs 1, 2, 3, 4, 5, 6, 7 et 8, chacune avec probabilité $\dfrac{1}{8}$.
\item On a $E(X_1)=\dfrac{1}{8} \times (1+2+3+4+5+6+7+8)=\dfrac{9}{2}$.
\item On a $E(S)=E(X_1+X_2+X_3)=E(X_1)+E(X_2)+E(X_3)=\dfrac{9}{2}+\dfrac{9}{2}+\dfrac{9}{2}=\dfrac{27}{2}$.
\item On a $S=24$ si et seulement si le numéro 8 a été tiré lors des trois tirages. Ainsi, $\pp(S=24)=\pp((8;8;8))=\dfrac{1}{512}$.
\item \begin{enumerate}
\item Les 10 tirages possibles sont $(8;8;8)$, $(8;8;7)$, $(8;7;8)$, $(7;8;8)$, $(8;7;7)$, $(7;8;7)$, $(7;7;8)$, $(8;8;6)$, $(8;6;8)$ et $(6;8;8)$.
\item La probabilité de gagner un lot est donc de $\dfrac{10}{512}$ soit $\dfrac{5}{256}$.
\end{enumerate}
\end{enumerate}

\end{solution}


\begin{exercise}

On lance trois pièces de monnaies et on regarde sur quels côtés elles tombent. 
\begin{enumerate}

\item On note $X$ le nombre de FACE obtenus. Construire le tableau résumant la loi de $X$.

\item On note $Y$ la variable aléatoire qui vaut 1 si les trois pièces tombent du même côté et 0 sinon. 
\begin{enumerate}
\item Quelle est la loi de $Y$ ? On précisera la valeur du ou des paramètres(s).
\item Que vaut l'espérance de $Y$ ?
\end{enumerate}
\item Le jeu consiste à miser deux euros. Si les trois pièces tombent sur les mêmes faces, on reprend sa mise et on remporte cinq euros supplémentaires. Sinon, on perd la mise. On note $Z$ la variable aléatoire qui détermine le gain algébrique du joueur.
\begin{enumerate}
\item Justifier que $Z=7Y-2$
\item En déduire l'espérance de $Z$. Ce jeu est-il équitable ?
\end{enumerate}
\end{enumerate}

\end{exercise}

\begin{solution}

$X$ suit une loi binomiale de paramètres $n=3$ et $p=0,5$. Le tableau résumant la loi de $X$ est le suivant :
\begin{center}
\begin{tabular}{|l|c|c|c|c|}
\hline
$k$ & $0$ & $1$ & $2$ & $3$ \\
\hline
$\mathbb{P}(Y=k)$ & $\dfrac{1}{8}$ & $\dfrac{3}{8}$ & $\dfrac{3}{8}$ & $\dfrac{1}{8}$ \\
\hline \end{tabular}\end{center}

$Y$ est une variable aléatoire de Bernoulli, de paramètre $\dfrac{1}{8}+\dfrac{1}{8}=\dfrac{1}{4}$. L'espérance de $Y$ vaut $\dfrac{1}{4}$.

On a alors $Z= 7Y-2$ : Si $Y$ vaut 0, on perd deux euros. Sinon, on en remporte 5. \\Ainsi, $E(Z)=7E(Y)-2=\dfrac{7}{4}-2=-\dfrac{1}{4}$. L'espérance est négative, ce jeu est donc au désavantage du joueur.
\end{solution}





\begin{exercise}
Une urne contient 100 jetons parmi lesquels 10 sont gagnants. Pour jouer à la loterie, un joueur doit payer 10 euros et tire au hasard et successivement deux jetons, en remettant entre temps le jeton tiré. Chaque jeton gagnant tiré lui rapporte 20 euros.

\begin{enumerate}
\item On note $X$ le nombre de jetons gagnants tirés. Quelle est la loi de $X$ ?
\item Que vaut l'espérance de $X$ ?
\item On note $Y$ le gain algébrique d'un joueur. Expliquer pourquoi $Y=20X-10$.
\item En déduire l'espérance de $Y$. Ce jeu est-il équitable ?
\end{enumerate}\newpage \end{exercise}

\begin{solution}On note $X$ le nombre de jetons gagnants tirés. $X$ suit une loi binomiale de paramètres 2 et $\dfrac{1}{10}$. L'espérance de $X$ vaut $2 \times \dfrac{1}{10}=0,2$

On note $Y$ le gain algébrique d'un joueur. Si $X$ représente le nombre de jetons gagnants tirés, le gain est de $20 X$. La participation au jeu étant de 10 euros, le gain $Y$ vaut $Y=20X-10$. \\On a donc $E(X)=20 E(X)-10=20 \times 0,2 -10 = -6$. Le jeu est au désavantage du joueur.\end{solution}







\begin{exercise}On considère les deux variables aléatoires $X$ et $Y$ indépendantes dont les lois sont données ci-dessous.

\renewcommand{\arraystretch}{2.2}
\hfill
\begin{tabular}{|l|c|c|c|}
\hline
$k$ & $2$& $1$ & $-1$  \\
\hline
$\mathbb{P}(X=k)$ & $\dfrac{1}{8}$ & $\dfrac{1}{4}$ & $\dfrac{5}{8}$ \\
\hline \end{tabular}\hfill
et
\hfill
\begin{tabular}{|l|c|c|c|}
\hline
$k$ & $1$& $2$ & $-2$ \\
\hline
$\mathbb{P}(Y=k)$ & $\dfrac{1}{2}$ & $\dfrac{1}{6}$ & $\dfrac{1}{3}$ \\
\hline \end{tabular}\hfill~

\begin{enumerate}
\item Donner l'espérance et la variance des variables aléatoires $X$ et $Y$.
\item On propose le jeu suivant : 8 boules sont dans une urnes. On mise un euro et on tire une de ces boules au hasard. 5 sont perdantes, 2 font gagner 2 euros et 1 fait gagner 3 euros. Quelle variable aléatoire permet de modéliser ce jeu ?
\item Le jeu est-il avantageux pour le joueur ?
\item Proposer une expérience aléatoire correspondant à la variable aléatoire $Y$.
\item On réalise deux fois le jeu correspondant à la variable $X$ et trois fois celui correspondant à la variable $Y$. On note $Z$ le gain algébrique de cette successions de jeu. Sans déterminer précisément la loi de $Z$, dire si ce jeu est avantageux pour le joueur ou non.
\end{enumerate}\end{exercise}

\begin{solution}On a  $E[X] = 2 \times \dfrac{1}{8} + 1 \times \dfrac{1}{4} + (-1) \times \dfrac{5}{8} = -\dfrac{1}{8}$ et
\[ V(X)= \dfrac{1}{2}\left(2- \left(-\dfrac{1}{8}\right)\right)^2+\dfrac{1}{4}\left(1-\left(-\dfrac{1}{8}\right)\right)^2+\dfrac{5}{8}\left(-1-\left(-\dfrac{1}{8}\right)\right)^2 = \dfrac{1563}{512}\simeq 3,05\]
Par ailleurs, $E[Y]=1 \times \dfrac{1}{2}+ 2 \times \dfrac{1}{6} + (-2) \times \dfrac{1}{3} = \dfrac{1}{6}$ et
\[ V(Y)= \dfrac{1}{3}\left(1- \dfrac{1}{6}\right)^2+\dfrac{1}{6}\left(2-\dfrac{1}{6}\right)^2+\dfrac{1}{2}\left(-2-\dfrac{1}{6}\right)^2 = \dfrac{89}{36}\simeq 2,47\]

La variable $X$ correspond au jeu proposé ici. L'espérance est négative : le jeu est désavantageux pour le joueur

Pour la loi de $Y$ : Dans une urne sont placées six boules. Une participation coûte deux euros. Deux boules sont perdantes, 3 rapportent 3 euros et 1 rapporte 4 euros.

On a $Z=2X+3Y$ et donc $E(Z)=2E(X)+ 3 E(Y)= -2\times \dfrac{1}{8} + 3 \times \dfrac{1}{6}=\dfrac{1}{4}>0$. Le jeu est avantageux pour le joueur.\end{solution}




\begin{exercise}Soit $X$ une variable aléatoire réelle suivant une loi binomiale. On suppose que $E(X)=6$ et $V(X)=4$. Retrouver les paramètres de la loi binomiale.\end{exercise}

\begin{solution}D'une part, $V(X)=(1-p)E(X)$. Ainsi, $p=\dfrac{1}{3}$. Par ailleurs, $E(X)=np=6$ d'où $n=18$.\end{solution}



\begin{exercise}
Soit $X$ une variable aléatoire suivant une loi binomiale $\mathcal{B}(10 ; 0,3)$ et $Y$ une variable aléatoire suivant une loi binomiale $\mathcal{B}(8;0,2)$. On suppose que $X$ et $Y$ sont des variables indépendantes.
\begin{enumerate}
\item Donner les espérances et variances de $X$ et $Y$. Donner leur écart-type arrondi au millième.
\item Donner l'espérance, la variance et l'écart-type arrondi au millième de la variable aléatoire $Z$ définie par $Z=2X-3Y$.
\end{enumerate}\end{exercise}

\begin{solution}On a $E(X)= 10 \times 0,3=3$, $V(X)= 10 \times 0,3 \times (1-0,3)= 2,1$ et $\sigma(X)=\sqrt{2,1}\simeq 1,449$ puis \\ $E(Y)= 8 \times 0,2 = 1,6$, $V(Y)=8 \times 0,2 \times (1-0,2)= 1,28$ et $\sigma(X)=\sqrt{1,28} \simeq 1,131$.

Par ailleurs, $E(Z)=2E(X)-3E(Y)=2 \times 3 - 3 \times 1,6 = 1,2$. De plus, $X$ et $Y$ étant indépendantes, on a $V(Z)=V(2X)+V(-3Y)=4V(X)+9V(Y)=19.92$ et $\sigma(Z)=\sqrt{19.92}\simeq 4,463$
\end{solution}




\begin{exercise}
On range $n$ objets dans une commode contenant $n$ tiroirs, chaque objet étant placé uniformément au hasard et indépendamment des autres objets dans l'un de ces tiroirs. L'objectif de cet exercice est de déterminer le nombre moyen de tiroirs vides à l'issue de cette expérience.
\vskip10pt
\textbf{Partie A : Étude de cas particuliers}

\begin{enumerate}
\item On suppose qu'il y a 2 tiroirs et 2 objets à ranger. 

\begin{enumerate}
\item On note $X$ la variable aléatoire qui compte le nombre de tiroirs vides à l'issue du rangement des 2 objets. Construire le tableau donnant la loi de $X$
\item En déduire l'espérance de $X$. 
\end{enumerate}
\item On suppose qu'il y a 3 tiroirs et 3 objets à ranger dans la commode. Un rangement peut alors être assimilé à un 3-uplet de $\{1;2;3\}$. Par exemple, le 3-uplet $(2;1;2)$ signifie que  le premier objet est rangé dans le tiroir 2, le deuxième objet dans le tiroir 1 et le troisième objet dans le tiroir 2.
\begin{enumerate}
\item Combien de rangements différents peut-on effectuer ?
\item Combien de rangements laissant le tiroir 1 vide peut-on effectuer ?
\item Pour $k \in \{ 1 ; 2 ; 3\}$ on note $X_k$ la variable aléatoire qui vaut $1$ si le tiroir $k$ est vide et 0 sinon. Quelle est la loi de $X_1$ ?
\item On considère la variable aléatoire $X=X_1+X_2+X_3$. Que vaut $E[X]$ ? Interpréter ce résultat dans le contexte de l'exercice.
\end{enumerate}
\end{enumerate}
\newpage
\textbf{Partie B : Cas général}

On dispose désormais de $n$ objets que l'on répartit uniformément au hasard et de manière indépendante dans les $n$ tiroirs.
\begin{enumerate}
\item Pour $k \in \{ 1 ; 2 ; \ldots ; n\}$ on note $Y_k$ la variable aléatoire qui vaut $1$ si le tiroir $k$ est vide et 0 sinon. \\ Montrer que l'espérance de $Y_1$ vaut $\dfrac{(n-1)^n}{n^n}$.
\item En déduire le nombre moyen de tiroirs vides à l'issue de cette expérience. 
\end{enumerate}\end{exercise}


\begin{solution}
\textbf{Partie A : Étude de cas particuliers}

\begin{enumerate}
\item 

\begin{enumerate}
\item La loi de $X$ est la suivante.

\begin{tabularx}{0.5\linewidth}{|X|X|X|}
\hline
$k$ & 0 & 1 \\
\hline
$\mathbb{P}(X=k)$ & $\dfrac{1}{2}$ & $\dfrac{1}{2}$\\
\hline
\end{tabularx}

La probabilité d'avoir un tiroir vide correspond au cas où les deux objets sont rangés dans le tiroir 1 (probabilité $\dfrac{1}{2} \times \dfrac{1}{2} = \dfrac{1}{4}$) ou dans le tiroir 2 (même probabilité), soit une probabilité total de $\dfrac{1}{2}$. La probabilité qu'il y ait un tiroir vide s'obtient par complément à 1.

\item $E[X]=\dfrac{1}{2} \times 0 + \dfrac{1}{2} \times 1 = \dfrac{1}{2}$.
\end{enumerate}
\item 
\begin{enumerate}
\item On peut effectuer $3 \times 3 \times 3 = 27$ rangements différents.
\item Si le tiroir 1 n'est pas utilisé, on peut effectuer $2 \times 2 \times 2 = 8$ rangements différents.
\item $X_1$ suit une loi de Bernoulli de paramètre $\dfrac{8}{27}$.
\item On considère la variable aléatoire $X=X_1+X_2+X_3$.  $X$ donne le nombre de tiroirs vides à l'issue de l'expérience. $E[X]=E[X_1]+E[X_2]+E[X_3] = 3 \times \dfrac{8}{27}=\dfrac{8}{9}$. Si l'on fait un grand nombre de fois cette expérience, la moyenne de tiroirs vide se rapprochera de $\dfrac{8}{9}$.
\end{enumerate}
\end{enumerate}

\textbf{Partie B : Cas général}

\begin{enumerate}
\item Il y a $n^n$ rangements dans les tiroirs possibles et $(n-1)^n$ rangements dans les tiroirs qui laissent le tiroir 1 vide. Ainsi, $\mathbb{P}(Y_1=1)=\dfrac{(n-1)^n}{n^n}$. Puisque $Y_1$ est une variable de Bernoulli, c'est également son espérance.
\item Le nombre moyen de tiroirs vides est l'espérance de $Y=Y_1+Y_2+\ldots + Y_n$. Les $Y_k$ ont tous pour espérance $\dfrac{(n-1)^n}{n^n}$. Ainsi, $E(X)=n \times \dfrac{(n-1)^n}{n^n}=\dfrac{(n-1)^n}{n^{n-1}}$.
\end{enumerate}

\end{solution}



\section*{Loi des grands nombres}


\begin{exercise}Soit $X$ une variable aléatoire d'espérance 4 et de variance 2. On considère un échantillon $(X_1, \ldots, X_n)$ de variables aléatoires indépendantes de même loi que $X$ et on note $M_n = \dfrac{1}{n} (X_1+X_2+\ldots + X_n)$.
\begin{enumerate}
\item Donner l'espérance et la variance de $M_n$.
\item Pour quelle valeur de $n$ la variance de $M_n$ est-elle inférieure à $10^{-4}$ ?
\end{enumerate}
\end{exercise}

\begin{solution}L'espérance de $M_n$ vaut 4 et sa variance $\dfrac{2}{n}$.
La variance de $M_n$ est inférieure à $10^{-4}$ dès que $n\geqslant 2 \times 10^{4}$.\end{solution}



\begin{exercise}Soit $X$ une variable aléatoire d'espérance 4 et de variance 1. 
\begin{enumerate}
\item Traduire l'inégalité $|X-4|\geqslant 2$ en terme d'intervalle.
\item Donner une minoration de $\mathbb{P}(|X-E(X)| \in ]2;6[)$.
\end{enumerate}
\end{exercise}

\begin{solution}On a $|X-4|\geqslant 2$ si et seulement si $X \in ]2;6[$.

D'après l'inégalité de Bienaymé-Tchebychev, $\mathbb{P}(|X-E(X)|\geqslant 2)\leqslant \dfrac{V(X)}{2^2} =\dfrac{1}{4}$.

Ainsi, $\mathbb{P}(|X-E(X)| \in ]2;6[) \geqslant \dfrac{3}{4}$.\end{solution}

\begin{exercise}Soit $X$ une variable aléatoire non constante. À l'aide de l'inégalité de Bienaymé-Thecbychev, minorer les probabilités $\pp(|X-E(X)|<2 \sigma(X))$ et $\pp(|X-E(X)|< 3 \sigma(X))$.\end{exercise}

\begin{solution}
D'après l'inégalité de Bienaymé-Tchebychev, on a $\pp(|X-E(X)|\geqslant 2\sigma (X)) \geqslant \dfrac{V(X)}{(2\sigma(X))^2}$. Or, $(2\sigma(X))^2=4\sigma(X)^2=4V(X)$.

Ainsi, $\pp(|X-E(X)|\geqslant 2\sigma (X)) \leqslant \dfrac{1}{4}$ et donc $\pp(|X-E(X)|< 2\sigma (X)) = 1- \pp(|X-E(X)|\geqslant 2\sigma (X)) \geqslant \dfrac{3}{4}$.

De même, on montrer que $\pp(|X-E(X)|< 3\sigma (X)) \geqslant \dfrac{8}{9}$.
\end{solution}


\begin{exercise}En 2018, le trafic moyen quotidien de véhicules sur le réseau autoroutier s'élevait à 24000 voitures, avec une variance de 6000. Majorer la probabilité que l'écart entre le nombre de véhicules en circulation lors d'une journée prise au hasard et la moyenne de véhicules recensés soit supérieure ou égal à 1000, puis à 100.\end{exercise}

\begin{solution} Notons $X$ la variable aléatoire qui donne le nombre de véhicules en circulation lors d'une journée prise au hasard. D'après l'énoncé, on a donc $E(X)=24000$ et $V(X)=6000$. D'après l'inégalité de Bienaymé-Tchebychev, la probabilité que l'écart entre le nombre de véhicules en circulation lors d'une journée prise au hasard et la moyenne de véhicules recensés soit supérieure ou égal à 1000 peut être majorée comme suit :
\[ \mathbb{P}(|X-E(X)| \geqslant 1000) \leqslant \dfrac{V(X)}{1000^2} = 0,006.\]
De la même manière,
\[ \mathbb{P}(|X-E(X)| \geqslant 100) \leqslant \dfrac{V(X)}{100^2} = 0,6.\]\end{solution}


\begin{exercise}[subtitle={(Métropole 2024, Jour 1 Bis)}]

Un client arrive à une station-service et se dirige vers une pompe. Il constate que deux voitures sont
devant lui, la première accédant à la pompe au moment de son arrivée.

On désigne par $T_1$, $T_2$, $T_3$ les variables aléatoires qui modélisent les temps passés en minute par chacun des trois clients, dans leur ordre d'arrivée, pour alimenter son véhicule entre l'instant où la pompe est disponible pour lui et celui où il la libère.

On suppose que $T_1$, $T_2$, $T_3$ sont des variables aléatoires indépendantes de même espérance égale à 6 et de même variance égale à 1.

On note $S$ la variable aléatoire correspondant au temps d'attente total passé à la station du troisième client entre son arrivée à la station et son départ de la pompe après avoir alimenté son véhicule.
\begin{enumerate}
\item Exprimer $S$ en fonction de $T_1$, $T_2$ et $T_3$.
\item \begin{enumerate}
\item Déterminer l'espérance de $S$ et interpréter ce résultat dans le contexte de l'exercice.
\item Quelle est la variance du temps d'attente total S de ce troisième client ?\end{enumerate}
\item Montrer que la probabilité que le troisième client passe un temps strictement compris entre 14 et 22 minutes à la station est supérieure ou égale à 0,81.
\end{enumerate}
\end{exercise}

\begin{solution}
On a $S=T_1+T_2+T_3$. Ainsi, $E(S)=E(T_1+T_2+T_3)=E(T_1)+E(T_2)+E(T_3)=18$. En moyenne, un automobiliste qui a deux personnes devant lui dans la file d'attente attendra 18 minutes avant de pouvoir repartir.

Par ailleurs, puisque $T_1$, $T_2$ et $T_3$ sont indépendantes, on a $V(S)=V(T_1)+V(T_2)V(T_3)=3$.

On cherche alors à minorer la probabilité $\pp(S \in ]14;22[)$, c'est-à-dire $\pp(S \in ]18-4;18+4[)$.

La probabilité recherchée est donc $\pp(|S-18|<4)$, c'est-à-dire $\pp(|S-E(S)|<4)$.

D'après l'inégalité de Bienaymé-Tchebychev, pour tout réel $\delta$ strictement positif, $\pp(|S-E(S)|\geqslant \delta) \leqslant \dfrac{V(S)}{\delta ^2}$. En prenant $\delta=4$, on a donc $\pp(|S-E(S)|\geqslant 4)\leqslant \dfrac{3}{16}$.

Ainsi, $\pp(|S-E(S)|<4)=1-\pp(|S-E(S)|\geqslant 4)\geqslant 1- \dfrac{3}{16}$. On a donc $\pp(|S-E(S)|<4)\geqslant \dfrac{13}{16}$. 

Or, $\dfrac{13}{16}=0,8125>0,81$. Finalement, la probabilité que le troisième client passe un temps strictement compris entre 14 et 22 minutes à la station est supérieure ou égale à 0,81.
\end{solution}

\begin{exercise}[subtitle={(Métropole 2024, jour 2)}]
La directrice d'une école souhaite réaliser une étude auprès des étudiants qui ont passé l'examen de fin d'étude.

On interroge au hasard dix étudiants. Les variables aléatoires $N_1$, $N_2$, ..., $N_{10}$ modélisent la note sur 20 obtenue à l'examen par chacun d'entre eux. On admet que ces variables sont indépendantes et suivent la même loi binomiale de paramètres $(20 ; 0,615)$.

Soit $S$ la variable définie par $S = N_1 + N_2 +··· + N_{10}$.
\begin{enumerate}
\item Calculer l'espérance $E(S)$ et la variance $V(S)$ de la variable aléatoire $S$.
\item On considère la variable aléatoire $M=\dfrac{S}{10}$.
\begin{enumerate}
\item Que modélise cette variable aléatoire M dans le contexte de l'exercice ?
\item Justifier que $E(M) = 12,3$ et $V (M) = 0,47355$.
\item À l'aide de l'inégalité de Bienaymé-Tchebychev, justifier l'affirmation ci-dessous. \\ « La probabilité que la moyenne des notes de dix étudiants pris au hasard soit
strictement comprise entre 10,3 et 14,3 est d'au moins 80% ».
\end{enumerate}
\end{enumerate}
\end{exercise}


\begin{exercise}On jette 3600 fois un dé équilibré à 6 faces,  numérotées de 1 à 6. On note $X$ la variable aléatoire qui compte le nombre de 1 obtenus.
\begin{enumerate}
\item Quelle est la loi de $X$ ? Quelle est son espérance ? sa variance ?
\item Minorer la probabilité que le nombre d'apparitions du numéro 1 soit compris entre 480 et 720.\end{enumerate}\end{exercise}

\begin{solution}$X$ suit une loi binomiale de paramètres $3600$ et $\dfrac{1}{6}$. \\Ainsi, $E[X]=3600 \times \dfrac{1}{6} = 600$ et $V(X)=3600 \times \dfrac{1}{6} \times \dfrac{5}{6}=500$.

Par ailleurs, $X \in ]480;720[$ correspond à $X-600 \in ]-120;120[$ c'est-à-dire $|X-600| < 120$. \\D'après l'inégalité de Bienaymé-Techebychev, 
$\mathbb{P}(|X-E[X]| \geqslant 120) \leqslant \dfrac{500}{120^2}=0,035$. \\Ainsi, $\mathbb{P}(|X-E[X]| < 120) \geqslant 1-0,035 = 0,965$.\end{solution}




\begin{exercise}Soit $X$ une variable aléatoire d'espérance 5 et de variance 2. On considère un échantillon $(X_1, \ldots, X_{100})$ de variables aléatoires indépendantes de même loi que $X$ et on note $M_n = \dfrac{1}{n} (X_1+X_2+\ldots + X_n)$.
\begin{enumerate}
\item Soit $\delta$ un réel strictement positif et $n$ un entier naturel non nul. \\ Écrire l'inégalité de concentration pour $M_n$.
\item En déduire l'entier $n$ à partir duquel on a $\mathbb{P}(|M_n-5|\geqslant 0,05 ) \leqslant 0,01$.
\end{enumerate}\end{exercise}

\begin{solution}L'inégalité de concentration pour $M_n$ s'écrit $\mathbb{P}(|M_n-E(X_1)|\geqslant \delta )\leqslant \dfrac{V(X)}{n\delta ^2}=\dfrac{2}{n\delta ^2}$.

En prenant $\delta = 0,05$, on obtient $\mathbb{P}(|M_n-5|\geqslant 0,05 ) \leqslant \dfrac{2}{0,0025n}$, qui est inférieur à 0,01 lorsque $n\geqslant 80000$.\end{solution}




\begin{exercise}Une compagnie aérienne exploite un avion ayant une capacité de 200 places. Pour ce vol, une analyse a montré que chaque passager à une probabilité $p=0,8$ de se présenter à l'embarquement. On suppose que les présences individuelles des passagers à l'embarquement sont indépendantes. La compagnie souhaite vendre davantage de billets que de places disponibles tout en limitant le risque de voir trop de personnes se présenter à l'embarquement.

Soit $n$ un entier strictement supérieur à 200, correspondant au nombre de billets vendus. On note $S_n$ le nombre de personnes se présentant à l'embarquement.
\begin{enumerate}
\item Quelle est la loi de $S_n$ ? Que valent son espérance et sa variance ?
\item On suppose que $n<250$. 
\begin{enumerate}
\item Justifier que si $S_n \geqslant 200$, alors $|S_n-0,8n| \geqslant 200-0,8n$.
\item En déduire que $\mathbb{P}(S_n \geqslant 200) \leqslant \dfrac{0,16n}{(200-0,8n)^2}$.
\item Combien de billets la compagnie peut-elle vendre tout en ayant une probabilité inférieure à 5\% que plus de 200 clients se présentent à l'embarquement ?
\end{enumerate}
\end{enumerate}
\newpage \end{exercise}

\begin{solution}$S_n$ suit une loi binomiale de paramètres $n$ et 0,8. Ainsi, $E[S_n]=0,8n$ et $V(S_n)=0,16n$.

Si $S_n \geqslant 200$, alors $S_n - 0,8n \geqslant 200-0,8n$. Puisque $n<250$, alors $200-0,8n >0$ et donc, par croissance de la fonction valeur absolue sur $\mathbb{R}_+$, $|S_n-0,8n| \geqslant 200-0,8n$.

Ainsi, $\mathbb{P}(S_n \geqslant 200) \leqslant \mathbb{P}(|S_n-0,8n| \geqslant 200-0,8n)$. Or, d'après l'inégalité de Bienaymé-Tchebychev, $\mathbb{P}(|S_n-0,8n| \geqslant 200-0,8n) \leqslant \dfrac{0,16n}{(200-0,8n)^2}$, d'où le résultat voulu.

On a $\dfrac{0,16n}{(200-0,8n)^2} \leqslant 0,05$ si et seulement si $0,032n^2-16,16n+2000 \geqslant 0$. Il s'agit d'un polynôme du second degré dont les racines valent environ $217,06$ et $287.94$. En particulier, si $n <250$, cette quantité est positive pour $n \leqslant 217$. La compagnie peut vendre 217 billets : elle aura alors moins de 5\% de chance que plus de 200 clients se présentent à l'embarquement.\end{solution}



\begin{exercise}Un joueur joue à la roulette en misant à chaque fois un euro sur une couleur (rouge ou noir). A chaque partie, il récupère sa mise et gagne un euro avec probabilité $\dfrac{18}{37}$. Sinon, il perd sa mise. Pour tout entier naturel $n$, on note $X_n$ son gain après $n$ parties et $Y_n$ le nombre de parties gagnés parmi les $n$ premières parties.
\begin{enumerate}
\item Quelle est la loi de $Y_n$ ? Que vaut son espérance et sa variance ?
\item Exprimer $X_n$ en fonction de $Y_n$ puis donner son espérance et sa variance.
\item A l'aide de l'inégalité de Bienaymé-Tchebychev, déterminer un entier $n$ à partir duquel la probabilité que le joueur ait perdu de l'argent après $n$ parties soit supérieure ou égale à 0,95.
\end{enumerate}\end{exercise}

\begin{solution}$Y_n$ suit une loi binomiale de paramètres $n$ et $\dfrac{18}{37}$. Son espérance vaut $\dfrac{18n}{37}$ et sa variance $\dfrac{306n}{1369}$.

On a par ailleurs $X_n=Y_n \times 1 + (n-Y_n) \times (-1)=2Y_n-n$.\\ Ainsi, $E[X_n]=2E[X_n]-n=-\dfrac{n}{37}$ et $V(X_n)=4V(Y_n)=\dfrac{1224n}{1369}$.

On cherche alors $P(X_n <0)$. Or, si $X_n \geqslant 0$, alors $X_n+\dfrac{n}{37} \geqslant  \dfrac{n}{37}$ et donc $\left| X_n + \dfrac{n}{37} \right| \geqslant \dfrac{n}{37}$. \\Ainsi, $P(X_n \geqslant 0) \leqslant \mathbb{P}\left(\left| X_n + \dfrac{n}{37} \right| \geqslant \dfrac{n}{37}\right)$.

Or, d'après l'inégalité de Bienaymé-Tchebychev, $\mathbb{P}\left(\left| X_n + \dfrac{n}{37} \right| \geqslant \dfrac{n}{37}\right)  \leqslant \dfrac{\frac{1224n}{1369}}{\frac{n^2}{37^2}}=\dfrac{1224}{n}$. 

Ainsi, en utilisant le complémentaire, $P(X_n <0) \geqslant 1-\dfrac{1224}{n}$. Or, $1-\dfrac{1224}{n} \geqslant 0.95$ si et seulement si $n \geqslant \dfrac{1224}{0,05}=22480$. 

Un joueur qui fait 22480 parties a plus de 95\% de chances de perdre de l'argent à l'issue de ces parties.\end{solution}


\section*{Exercices de synthèse}

\begin{exercise}[subtitle={(Sujet zéro 2024)}]

Dans un examen, une épreuve notée sur dix points est constituée de deux exercices : le premier est noté sur deux points, le deuxième sur huit points.

\paragraph{Partie I}

Le premier exercice est constitué de deux questions Q1 et Q2.

Chaque question est notée sur un point. Une réponse correcte rapporte un point ; une réponse incorrecte, incomplète ou une absence de réponse rapporte zéro point.

On considère que :
\begin{itemize}
\item Un candidat pris au hasard a une probabilité 0,8 de répondre correctement à la question Q1.
\item Si le candidat répond correctement à Q1, il a une probabilité 0,6 de répondre correctement à Q2 ; s'il ne répond pas correctement à Q1, il a une probabilité 0,1 de
répondre correctement à Q2.\end{itemize}

On prend un candidat au hasard et on note :
\begin{itemize}
\item $A$ l'évènement : « le candidat répond correctement à la question Q1 » ;
\item $B$ l'évènement : « le candidat répond correctement à la question Q2 ».
\end{itemize}

On note $\overline{A}$ et $\overline{B}$ les évènements contraires de $A$ et de $B$.

\begin{enumerate}
\item Recopier et compléter les pointillés de l'arbre pondéré ci-dessous.

\tikzstyle{level 1}=[level distance=3.5cm, sibling distance=3cm]
\tikzstyle{level 2}=[level distance=3.5cm, sibling distance=1.5cm]
\tikzstyle{level 3}=[level distance=3.5cm, sibling distance=0.3cm]

% Define styles for bags and leafs
\tikzstyle{bag} = []
\tikzstyle{end} = [circle, minimum width=3pt,fill, inner sep=0pt]


\begin{center}
\begin{tikzpicture}[scale=0.7,grow=right,sloped]
\node[bag] { }
	child {
        node[bag] {$\overline{A}$} 
        child {
                node[bag] {$\overline{B}$}
                edge from parent node[above] {$\ldots$}
            }
        child {
               node[bag] {$B$}
               edge from parent node[above] {$\ldots$}
            }
            edge from parent node[above] {$\ldots$}
    }
	child {
        node[bag] {$A$} 
        child {
                node[bag] {$\overline{B}$}
                edge from parent node[above] {$\ldots$}
            }
        child {
               node[bag] {$B$}
               edge from parent node[above] {$\ldots$}
            }
            edge from parent node[above] {$\ldots$}
    };
\end{tikzpicture}
\end{center}

\item Calculer la probabilité que le candidat réponde correctement aux deux questions Q1 et Q2.
\item Calculer la probabilité que le candidat réponde correctement à la question Q2.\end{enumerate}
On note :
\begin{itemize}
\item $X_1$ la variable aléatoire qui, à un candidat, associe sa note à la question Q1 ;
\item $X_2$ la variable aléatoire qui, à un candidat, associe sa note à la question Q2 ;
\item $X$ la variable aléatoire qui, à un candidat, associe sa note à l'exercice, c'est-à-dire $X = X_1 + X_2$.
\end{itemize}

\begin{enumerate}
\setcounter{enumi}{3}
\item Déterminer l'espérance de $X_1$ et de $X_2$. En déduire l'espérance de $X$.\\ Donner une interprétation de l'espérance de $X$ dans le contexte de l'exercice.
\item On souhaite déterminer la variance de $X$.
\begin{enumerate}
\item Déterminer $\pp(X = 0)$ et $\pp(X = 2)$. En déduire $\pp(X = 1)$.
\item Montrer que la variance de $X$ vaut 0,57.
\item A-t-on $V(X) = V(X_1)+V(X_2)$ ? Est-ce surprenant ?
\end{enumerate} 
\end{enumerate}

 \paragraph{Partie II}
 
Le deuxième exercice est constitué de huit questions indépendantes.

Chaque question est notée sur un point. Une réponse correcte rapporte un point ; une réponse incorrecte et une absence de réponse rapporte zéro point.

Les huit questions sont de même difficulté : pour chacune des questions, un candidat a une probabilité $\dfrac{3}{4}$ de répondre correctement, indépendamment des autres questions.

On note $Y$ la variable aléatoire qui, à un candidat, associe sa note au deuxième exercice, c'est-à-dire le nombre de bonnes réponses.

\begin{enumerate}
\item Justifier que $Y$ suit une loi binomiale dont on précisera les paramètres.
\item Donner la valeur exacte de $\pp(Y = 8)$.
\item Donner l'espérance et la variance de $Y$.
\end{enumerate}

\paragraph{Partie III}

On suppose que les deux variables aléatoires $X$ et $Y$ sont indépendantes. On note la variable aléatoire qui, à un candidat, associe sa note totale à l'examen : $Z = X +Y$.

\begin{enumerate}
\item Calculer l'espérance et la variance de $Z$.
\item Soit $n$ un nombre entier strictement positif.\\
Pour $i$ entier variant de 1 à $n$, on note $Z_i$ la variable aléatoire qui, à un échantillon de $n$ élèves, associe la note de l'élève numéro $i$ à l'examen.

On admet que les variables aléatoires $Z_1$, $Z_2$, ..., $Z_n$ sont identiques à $Z$ et indépendantes.

On note $M_n$  la variable aléatoire qui, à un échantillon de $n$ élèves, associe la moyenne de leurs $n$ notes, c'est-à-dire : $M_n=\dfrac{Z_1+Z_2+\ldots+Z_n}{n}$.
\vskip5pt
\begin{enumerate}
\item Quelle est l'espérance de $M_n$ ?
\item Quelles sont les valeurs de $n$ telles que l'écart type de $M_n$ soit inférieur ou égal à 0,5 ?
\item Pour les valeurs trouvées en \textbf{b.}, montrer que la probabilité que $6,3 < M_n < 8,3$ est supérieure ou égale à 0,75 .
\end{enumerate}
\end{enumerate}

\end{exercise}

\begin{solution}
\paragraph{Partie I}

\begin{enumerate}
\item On complète l'arbre de probabilités modélisant la situation

\tikzstyle{level 1}=[level distance=3.5cm, sibling distance=3cm]
\tikzstyle{level 2}=[level distance=3.5cm, sibling distance=1.5cm]
\tikzstyle{level 3}=[level distance=3.5cm, sibling distance=0.3cm]

% Define styles for bags and leafs
\tikzstyle{bag} = []
\tikzstyle{end} = [circle, minimum width=3pt,fill, inner sep=0pt]


\begin{center}
\begin{tikzpicture}[scale=0.7,grow=right,sloped]
\node[bag] { }
	child {
        node[bag] {$\overline{A}$} 
        child {
                node[bag] {$\overline{B}$}
                edge from parent node[above] {$0,9$}
            }
        child {
               node[bag] {$B$}
               edge from parent node[above] {$0,1$}
            }
            edge from parent node[above] {$0,2$}
    }
	child {
        node[bag] {$A$} 
        child {
                node[bag] {$\overline{B}$}
                edge from parent node[above] {$0,4$}
            }
        child {
               node[bag] {$B$}
               edge from parent node[above] {$0,6$}
            }
            edge from parent node[above] {$0,8$}
    };
\end{tikzpicture}
\end{center}

\item On a $\pp(A \cap B)= 0,8 \times 0,6=0,48$.

\item $(A;\overline{A})$ est un système complet d'événements. D'après la formule des probabilités totales,
\[ \pp(B)=\pp(A \cap B)+\pp(\overline{A} \cap B)=0,6 \times 0,8 +0,2 \times 0,1=0,5.\]
\item $X_1$ prend les valeurs 0 et 1, elle suit donc une loi de Bernoulli. Son paramètre (et donc son espérance) vaut 0,8.

De même, $X_2$ suit une loi de Bernoulli de paramètre 0,5. On a donc $E(X_2)=0,5$.

Ainsi, $E(X)=E(X_1+X_2)=E(X_1)+E(X_2)=0,8+0,5=1,3$. En moyenne, un élève obtient une note de 1,3 au premier exercice.

\item On a $\pp(X=0)=\pp(\overline{A}\cap \overline{B})=0,2 \times 0,9=0,18$ et $\pp(X=2)=\pp(A\cap B)=0,48$. \\Ainsi, $\pp(X=1)=1-(0,18+0,48)=0,34$.

\item On a $E(X^2)=0^2 \times 0,18 + 1^2 \times 0,34 + 2^2 \times 0,48 = 2,26$. 

D'après la formule de Koenig-Huygens, $V(X)=E(X^2)-E(X)^2=2,26-1,3^2=0,57$.

\item On a $V(X_1)=0,8 \times(1-0,8)=0,16$ et $V(X_2)=0,5 \times (1-0,5)=0,25$.\\ Ainsi, $V(X_1)+V(X_2)=0,41\neq V(X)$.

Ce résultat n'est pas surprenant puisque les variables aléatoires $X_1$ et $X_2$ ne sont pas indépendantes.

\end{enumerate}

\paragraph{Partie II}

\begin{enumerate}
\item $Y$ compte le nombre de succès (la réponse donnée est exacte) d'une répétitions de 8 épreuves identiques et indépendantes. La probabilité de succès est de $\dfrac{3}{4}$. Ainsi, $Y$ suit une loi binomiale $\mathcal{B}\left(8;\dfrac{3}{4}\right)$.
\item On a $\pp(Y=8)=\dbinom{8}{8} \times \left(\dfrac{3}{4}\right)^8\times\left(1-\dfrac{3}{4}\right)^{8-8}=\left(\dfrac{3}{4}\right)^8$.
\item On a $E(Y)=8 \times \dfrac{3}{4}=6$ et $V(Y)=8 \times \dfrac{3}{4} \times \left(1-\dfrac{3}{4}\right)=\dfrac{6}{4}=\dfrac{3}{2}$.
\end{enumerate}

\paragraph{Partie III}

\begin{enumerate}
\item On a $E[Z]=E[X+Y]=E[X]+E[Y] = 1,3+6=7,3$.\\ De plus, puisque $X$ et $Y$ sont indépendantes, on a $V(Z)=V(X+Y)=V(X)+V(Y)=0,57+1,5=2,07$.
\item \begin{enumerate}
\item On a $E[M_n]=E[Z_1]=7,3$.
\item On a $V(M_n)=\dfrac{V(Z_1)}{n}=\dfrac{2,07}{n}$ et donc $\sigma(M_n)=\sqrt{\dfrac{2,07}{n}}$.

Ainsi, $\sigma(M_n)\leqslant 0,5 \Leftrightarrow \sqrt{\dfrac{2,07}{n}} \leqslant 0,5 \Leftrightarrow \dfrac{\sqrt{2,07}}{0,5} \leqslant \sqrt{n}$. Par croissance de la fonction $x\mapsto x^2$ sur $[0;+\infty[$, on a alors $n \geqslant \dfrac{2,07}{0,25}$. Or, $\dfrac{2,07}{0,25}=8,28$. L'entier recherché est donc $n=9$.
\item On a $6,3 < M_n <8,3$ si et seulement si $M_n \in ]6,3 ; 8,3[$ soit $M_n \in ]7,3-1 ; 7,3+1[$ c'est-à-dire \\$|M_n-7,3|<1$.

Or, $E(M_n)=7,3$. D'après l'inégalité de Bienaymé-Tchebychev, on a $\pp(|M_n-7,3| \geqslant 1) \leqslant \dfrac{V(M_n)}{1^2}$ et donc $\pp(|M_n-7,3| \geqslant 1) \leqslant \dfrac{2,07}{n}$.

Ainsi, $\pp(|M_n-7,3|<1)=1-\pp(|M_n-7,3| \geqslant 1)\geqslant 1- \dfrac{2,07}{n}$. 

D'après la question précédente, on a $n\geqslant 9$. 

Ainsi, $\dfrac{2,07}{n} \leqslant \dfrac{2,07}{9}$ et $1-\dfrac{2,07}{n}\geqslant 1-\dfrac{2,07}{9}$. Or, $1-\dfrac{2,07}{9}=0,77$.

Ainsi, on a bien $\pp(|M_n-7,3|<1) \geqslant 0,77 \geqslant 0,75$.
\end{enumerate}
\end{enumerate}
\newpage
\end{solution}


\begin{exercise}[subtitle={(Métropole 2024, Jour 1)}]
Une agence de marketing a étudié la satisfaction des clients concernant le service clientèle à l'occasion de l'achat d'un téléviseur. Ces achats ont été réalisés soit sur internet, soit dans une chaîne de magasins d'électroménager, soit dans une enseigne de grandes surfaces.

Les achats sur internet représentent 60\% des ventes, les achats en magasin d'électroménager 30\% des ventes et ceux en grandes surfaces 10\% des ventes.

Une enquête montre que la proportion des clients satisfaits du service clientèle est de :
\begin{itemize}
\item 75\% pour les clients sur internet ;
\item 90\% pour les clients en magasin d'électroménager ;
\item 80\% pour les clients en grande surface.
\end{itemize}
On choisit au hasard un client ayant acheté le modèle de téléviseur concerné. On définit les évènements suivants.
\begin{itemize}
\item $I$ : « le client a effectué son achat sur internet » ;
\item $M$ : « le client a effectué son achat en magasin d'électroménager » ;
\item $G$ : « le client a effectué son achat en grande surface » ;
\item $S$ : « le client est satisfait du service clientèle ».
\end{itemize}

Si $A$ est un évènement quelconque, on notera $\overline{A}$ son évènement contraire et $P(A)$ sa probabilité.
\begin{enumerate}
\item Reproduire et compléter l'arbre ci-dessous.

\tikzstyle{level 1}=[level distance=3.5cm, sibling distance=2cm]
\tikzstyle{level 2}=[level distance=3.5cm, sibling distance=1cm]
\tikzstyle{level 3}=[level distance=3.5cm, sibling distance=0.3cm]

% Define styles for bags and leafs
\tikzstyle{bag} = []
\tikzstyle{end} = [circle, minimum width=3pt,fill, inner sep=0pt]

\begin{center}
\begin{tikzpicture}[scale=0.7,grow=right,sloped]
\node[bag] { }
	child {
        node[bag] {$G$} 
        child {
                node[bag] {$\overline{S}$}
                edge from parent node[above] {$\ldots$}
            }
        child {
               node[bag] {$S$}
               edge from parent node[above] {$\ldots$}
            }
            edge from parent node[above] {$\ldots$}
    }
	child {
        node[bag] {$M$} 
        child {
                node[bag] {$\overline{S}$}
                edge from parent node[above] {$\ldots$}
            }
        child {
               node[bag] {$S$}
               edge from parent node[above] {$\ldots$}
            }
            edge from parent node[above] {$\ldots$}
    }
    child {
        node[bag] {$I$} 
        child {
                node[bag] {$\overline{S}$}
                edge from parent node[above] {$\ldots$}
            }
        child {
               node[bag] {$S$}
               edge from parent node[above] {$\ldots$}
            }
            edge from parent node[above] {$\ldots$}
    };
\end{tikzpicture}
\end{center}

\item Calculer la probabilité que le client ait réalisé son achat sur internet et soit satisfait
du service clientèle.
\vskip5pt
\item Démontrer que $P(S) = 0,8$.
\vskip5pt
\item Un client est satisfait du service clientèle.\\ Quelle est la probabilité qu'il ait effectué son achat sur internet ? On donnera un résultat arrondi à $10^{-3}$ près.
\vskip5pt
\item Pour réaliser l'étude, l'agence doit contacter chaque jour 30 clients parmi les acheteurs du téléviseur. On suppose que le nombre de clients est suffisamment important
pour assimiler le choix des 30 clients à un tirage avec remise. On note $X$ la variable aléatoire qui, à chaque échantillon de 30 clients, associe le nombre de clients satisfaits du service clientèle.
\begin{enumerate}
\item Justifier que $X$ suit une loi binomiale dont on précisera les paramètres.
\item Déterminer la probabilité, arrondie à $10^{-3}$ près, qu'au moins 25 clients soient satisfaits dans un échantillon de 30 clients contactés sur une même journée.\end{enumerate}
\vskip5pt
\item En résolvant une inéquation, déterminer la taille minimale de l'échantillon de clients à contacter pour que la probabilité qu'au moins l'un d'entre eux ne soit pas satisfait soit supérieure à 0,99.
\vskip5pt
\item Dans les deux questions \textbf{a.} et \textbf{b.} qui suivent, on ne s'intéresse qu'aux seuls achats sur
internet.
\vskip5pt
Lorsqu'une commande de téléviseur est passée par un client, on considère que le
temps de livraison du téléviseur est modélisé par une variable aléatoire $T$ égale à la
somme de deux variables aléatoires $T_1$ et $T_2$.
\vskip5pt
La variable aléatoire $T_1$ modélise le nombre entier de jours pour l'acheminement du téléviseur depuis un entrepôt de stockage vers une plateforme de distribution. La
variable aléatoire $T_2$ modélise le nombre entier de jours pour l'acheminement du téléviseur depuis cette plateforme jusqu'au domicile du client.
\vskip5pt
On admet que les variables aléatoires $T_1$ et $T_2$ sont indépendantes, et on donne :
\begin{itemize}
\item L'espérance $ (T_1) = 4$ et la variance $V (T_1) = 2$ ;
\item L'espérance $E (T_2) = 3$ et la variance $V (T_2) = 1$.
\end{itemize}
\vskip5pt
\begin{enumerate}
\item Déterminer l'espérance $E(T )$ et la variance $V (T )$ de la variable aléatoire $T$.
\item Un client passe une commande de téléviseur sur internet. Justifier que la probabilité qu'il reçoive son téléviseur entre 5 et 9 jours après sa commande est
supérieure ou égale à $\dfrac{2}{3}$.
\end{enumerate}
\end{enumerate}


\end{exercise}

\begin{solution}\hspace{0pt}
\begin{enumerate}
\item On complète l'arbre comme suit.

\tikzstyle{level 1}=[level distance=3.5cm, sibling distance=2cm]
\tikzstyle{level 2}=[level distance=3.5cm, sibling distance=1cm]
\tikzstyle{level 3}=[level distance=3.5cm, sibling distance=0.3cm]

% Define styles for bags and leafs
\tikzstyle{bag} = []
\tikzstyle{end} = [circle, minimum width=3pt,fill, inner sep=0pt]

\begin{center}
\begin{tikzpicture}[scale=0.7,grow=right,sloped]
\node[bag] { }
	child {
        node[bag] {$G$} 
        child {
                node[bag] {$\overline{S}$}
                edge from parent node[below] {$0,2$}
            }
        child {
               node[bag] {$S$}
               edge from parent node[above] {$0,8$}
            }
            edge from parent node[above] {$0,1$}
    }
	child {
        node[bag] {$M$} 
        child {
                node[bag] {$\overline{S}$}
                edge from parent node[below] {$0,1$}
            }
        child {
               node[bag] {$S$}
               edge from parent node[above] {$0,9$}
            }
            edge from parent node[above] {$0,3$}
    }
    child {
        node[bag] {$I$} 
        child {
                node[bag] {$\overline{S}$}
                edge from parent node[below] {$0,25$}
            }
        child {
               node[bag] {$S$}
               edge from parent node[above] {$0,75$}
            }
            edge from parent node[above] {$0,6$}
    };
\end{tikzpicture}
\end{center}

\item On a $P(I \cap S)=0,6 \times 0,75 = 0,45$.

\item $(I;M;G)$ forme un système complet d'événements. D'après la formule des probabilités totales,
\[P(S)=P(I\cap S)+P(M\cap S)+P(G\cap S)=0,45+0,3 \times 0,9+0,1 \times 0,8 = 0,8.\]
\item On a $P_S(I)=\dfrac{P(S \cap I)}{P(S)}=\dfrac{0,45}{0,8}\simeq 0,563$.
\item \begin{enumerate}
\item Puisque les tirages sont assimilés à des tirages avec remise, $X$ compte le nombre de succès (le client est satisfait) d'une répétition de 30 expériences de Bernoulli identiques et indépendantes. $X$ suit une loi binomiale $\mathcal{B}\left(30;0,8\right)$.
\item D'après la calculatrice, on a $P(X \geqslant 25) \simeq 0,428$.
\end{enumerate}
\item Notons $Y$ le nombre de clients non satisfaits lorsque l'on interroge $n$ clients.\\ $Y$ suit une loi binomiale $\mathcal{B}(n;0,2)$.

On a alors $P(Y \geqslant 1)=1-P(Y <1)=1-P(Y=0)$ puisque $Y$ est à valeurs entières.\\ Or, $P(Y=0)=\dbinom{8}{0} \times 0,2 ^0 \times (1-0,2)^n=0,8^n$.

Ainsi, $P(Y \geqslant 1) \geqslant 0,99 \Leftrightarrow  1-0,8^n \geqslant 0,8^n \Leftrightarrow -0,8^n \geqslant -0,01 \Leftrightarrow 0,8^n \leqslant 0,01$.

Par croissance du logarithme népérien sur $]0;+\infty[$, on a alors $\ln(0,8^n) \leqslant \ln(0,01)$ soit $n\ln(0,8) \leqslant \ln(0,01)$ et donc $n \geqslant \dfrac{\ln(0,01)}{\ln(0,8)}$. Or, $\dfrac{\ln(0,01)}{\ln(0,8)}\simeq 20,6$. L'entier recherché est donc 21.
\item \begin{enumerate}
\item On a $E(T)=E(T_1+T_2)=E(T_1)+E(T_2)=4+3=7$.\\ De plus, puisque $T_1$ et $T_2$ sont indépendantes, $V(T)=V(T_1+T_2)=V(T_1)+V(T_2)=2+1=3$.
\item On cherche $P(5 \leqslant T \leqslant 9)$.\\ Puisque $T$ est à valeurs entières, cette probabilité est égale à $P(4<T<10)$. Or, on a $4<T<10$ si et seulement si $7-3<T<7+3$ c'est-à-dire $|T-7|<3$ et donc $|T-E(T)|<3$.

D'après l'inégalité de Bienaymé-Tchebychev, on a $P(|T-E(T)| \geqslant 3) \leqslant \dfrac{V(T)}{3^2}$.

On obtient donc $P(|T-E(T)| \geqslant 3) \leqslant \dfrac{1}{3}$.

Ainsi, $P(|T-E(T)|<3)=1-P(|T-E(T)| \geqslant 3) \geqslant 1- \dfrac{1}{3}$ et finalement, $P(|T-E(T)|<3) \geqslant \dfrac{2}{3}$.
\end{enumerate}
\end{enumerate}
\end{solution}




\chapter{Corrigés}


\printsolutions[headings={false} ]

\end{document}